\documentclass{article}

\PassOptionsToPackage{authoryear,round}{natbib}
\usepackage{neurips_2026}

\usepackage[utf8]{inputenc}
\usepackage[T1]{fontenc}
\usepackage{amsmath,amssymb,amsthm,mathtools}
\usepackage{amsfonts}
\usepackage{bbm}
\usepackage{algorithm}
\usepackage[noend]{algpseudocode}
\usepackage{booktabs}
\usepackage{enumitem}
\usepackage{xcolor}
\usepackage[colorlinks=true,linkcolor=blue,citecolor=blue,urlcolor=blue]{hyperref}
\usepackage[nameinlink,noabbrev]{cleveref}

\theoremstyle{plain}
\newtheorem{theorem}{Theorem}
\newtheorem{lemma}{Lemma}
\newtheorem{proposition}{Proposition}
\newtheorem{corollary}{Corollary}

\theoremstyle{definition}
\newtheorem{remark}{Remark}
\newtheorem{assumption}{Assumption}

\crefname{assumption}{Assumption}{Assumptions}
\Crefname{assumption}{Assumption}{Assumptions}
\crefname{lemma}{Lemma}{Lemmas}
\Crefname{lemma}{Lemma}{Lemmas}

\newcommand{\E}{\mathbb{E}}
\newcommand{\Var}{\mathrm{Var}}
\newcommand{\Prob}{\mathbb{P}}
\newcommand{\R}{\mathbb{R}}
\newcommand{\1}{\mathbbm{1}}
\newcommand{\Pzero}{P_0}
\newcommand{\psihat}{\widehat{\psi}_{\mathrm{cal}}}
\newcommand{\mhat}{m}
\newcommand{\fhat}{f_n}
\newcommand{\what}{\widehat w}
\newcommand{\wzero}{w_0}
\newcommand{\rhozero}{\rho_0}
\newcommand{\rhon}{\rho_n}
\newcommand{\mtilde}{m_n^\star}
\newcommand{\muzero}{\mu_0}
\newcommand{\Pl}{\mathbb{P}^{L}_n}
\newcommand{\Pu}{\mathbb{P}^{U}_N}
\newcommand{\Px}{P_{0,X}}
\newcommand{\M}{\mathcal M}
\newcommand{\Cov}{\operatorname{Cov}}


\usepackage{fvextra}

\DefineVerbatimEnvironment{codeblock}{Verbatim}{
  breaklines=true,
  breakanywhere=true,
  fontsize=\small
}

\title{Prediction-Powered Inference via Calibration}

\author{Lars van der Laan, Mark van der Laan}
\date{\today}

\begin{document}

\maketitle
\begin{abstract}
We study semisupervised mean estimation with one labeled sample and one unlabeled sample, where a black-box prediction model is available but may be biased. A standard approach in this setting is augmented inverse probability weighted (AIPW) estimation \citep{robins1995analysis}, which leverages the prediction model to improve inference while remaining valid even when that model is inaccurate. We propose a class of two-step regression-based estimators that first recalibrate the prediction model on the labeled sample and then average the recalibrated predictions over the pooled covariate sample. The resulting estimator retains a simple imputation-based plug-in form, but also admits an AIPW representation in which calibration removes the augmentation term. We implement calibration using isotonic regression, a lightweight, tuning-free method that learns the best monotone transformation of the prediction score. We show that the resulting estimator is root-\(n\) consistent and asymptotically normal under weak conditions. It is statistically efficient among estimators based on the calibrated score, and fully efficient when the calibrated limit equals the true regression function. We also clarify the relationship with recent prediction-powered inference (PPI) methods: PPI and PPI++ are special cases of AIPW, while PPI++ is also algebraically equivalent to linear calibration and classical linear regression adjustment based on prognostic scores. These linear methods are often less efficient than the isotonic-calibrated plug-in estimator we propose.
\end{abstract}

\section{Introduction}
We study semisupervised mean estimation with one labeled dataset \(\{(X_i, Y_i)\}_{i=1}^n\) and one unlabeled dataset \(\{\widetilde X_i\}_{i=1}^N\). This is a classical missing-outcome problem in semiparametric statistics, missing-data theory, and debiased machine learning \citep{rubin1976,robinsCausal,robins1995analysis,bang2005doubly,vanderlaanunified,vanderLaanRose2011,DoubleML}. Efficient statistical inference with black-box machine learning is possible using standard debiased machine-learning techniques such as augmented inverse-probability weighting, targeted minimum loss estimation, and double machine learning \citep{li2023efficient,graham2024towards}. Recent work in machine learning has referred to this setting as prediction-powered inference, especially in applications with black-box predictors and abundant unlabeled data \citep{ppi2023,crossppi2024}.
 

In the standard missing-at-random (MAR) setup, the labeled data are used to debias a naive plug-in estimator. If \(m(X)\) is a black-box prediction score, the naive plug-in estimator averages \(m(X)\) over the pooled covariate sample, while the labeled sample is used to estimate the residual correction term \(\mathbb{E}\{Y-m(X)\}\). For example, the standard augmented inverse probability weighted (AIPW) estimator \citep{robins1995analysis} is \(\frac{1}{N+n}\{\sum_{i=1}^n m(X_i)+\sum_{i=1}^N m(\widetilde X_i)\}+\frac{1}{n}\sum_{i=1}^n \{Y_i-m(X_i)\}\), whereas the PPI estimator \citep{ppi2023}, given by \(\frac{1}{N}\sum_{i=1}^N m(\widetilde X_i)+\frac{1}{n}\sum_{i=1}^n \{Y_i-m(X_i)\}\), is a typically less efficient variant that discards covariate information from the labeled sample. These estimators are effective for bias reduction, but they correct only the mean bias of the score.


Calibration provides a stronger and widely studied notion of predictive reliability \citep{zadrozny2001obtaining,zadrozny2002transforming,niculescu2005predicting}. Informally, a calibrated score is one whose numerical values accurately reflect the observed outcomes it predicts. Rather than correcting only the average bias \(\E[Y-m(X)]\), one can use the labeled sample to construct a calibrated score \(m_n^\star(X)\) such that, ideally, \(\E[Y \mid m_n^\star(X)]=m_n^\star(X)\). This stronger property immediately implies \(\E[m_n^\star(X)]=\E[Y]\), so the resulting plug-in estimator is automatically calibrated for the mean. This is attractive in practice because a useful prediction model may already be available, while a relatively small labeled sample can often recalibrate its output without retraining the underlying model \citep{roth2022uncertain}. Such recalibration can be implemented with classical post-processing methods such as Platt scaling \citep{platt1999probabilistic}, histogram binning \citep{zadrozny2001obtaining}, or isotonic regression \citep{zadrozny2002transforming}. This is especially relevant for modern machine-learning predictors, such as neural networks, which are often poorly calibrated \citep{bella2010calibration,guo2017calibration,wang2023calibration}.

 

Our proposal is therefore simple: fit a black-box regression model, calibrate it on the labeled sample, and then average the calibrated predictions over the pooled covariate sample. The resulting estimator is still just ``predict, calibrate, and average,'' but our main result shows that it also admits a debiased interpretation. In particular, it is exactly an AIPW estimator constructed from the calibrated scores, for which calibration eliminates the augmentation term. Thus, calibration turns a simple plug-in rule into a debiased semiparametric estimator \citep{robins1995analysis,bang2005doubly,vanderLaanRose2011,DoubleML} with strong efficiency guarantees. 


\paragraph{Contributions.}
We make three main contributions. First, in Section~\ref{sec:calppi}, we propose a calibrated plug-in estimator for semisupervised mean estimation under the joint two-sample model. The estimator combines a black-box regression fit with a calibration step on the labeled sample and then averages the calibrated predictions over the pooled covariate distribution. It can be implemented using classical calibration methods such as linear calibration, histogram binning, and isotonic regression; see, in particular, Sections~\ref{sec:iso} and \ref{sec:linear}. Among these, isotonic calibration is a lightweight, tuning-free method that learns the best monotone transformation of a score \citep{zadrozny2002transforming,niculescu2005predicting}. Our estimator is an instance of calibrated debiased machine learning \citep{van2024automatic}. We provide synthetic and real-data experiments in Section~\ref{sec:exp} that demonstrate the effectiveness of our approach. \texttt{Python} code is provided in Appendix~\ref{appendix:code}.

Second, in Section~\ref{sec:calppi} we show that calibration makes the predictions satisfy empirical score equations and that the resulting estimator admits an exact AIPW representation. Thus, isotonic calibration automatically debiases the seemingly naive imputation-based plug-in estimator, enabling root-\(n\) inference. We also study linear calibration in Section~\ref{sec:linear} and show that the AIPW estimators termed PPI and PPI++ are algebraically equivalent to linear regression adjustment using \(m(X)\) as a prognostic score, a classical approach to improving precision through flexible covariate adjustment in small-sample settings with randomized treatment or missing-at-random outcomes \citep{hansen2008prognostic,rosenblum2009using,hojbjerre2025powering}.  

Third, in Section~\ref{sec:theory} we derive asymptotic normality, valid inference, and efficiency results for our isotonic-calibrated estimator under the joint two-sample model. Confidence intervals can be constructed using standard Wald-type methods or the bootstrap. We show that the estimator is efficient in the reduced model where the best monotone transformation of the prediction score serves as a one-dimensional summary of the features, and is therefore more efficient than the AIPW estimator based on residuals from the uncalibrated score. If this monotone transformation coincides with the true outcome regression, then the estimator is nonparametrically efficient.  We review the two-sample semisupervised inference problem and its corresponding efficiency theory in Section \ref{sec:setup}.

 


\subsection{Related work}
\textbf{Semiparametric statistics and missing data.} This paper builds on foundational ideas from the missing-data and semiparametric literatures. Rubin formalized missing at random as a key identification condition for recovering full-data targets from partially observed outcomes \citep{rubin1976}. Semiparametric efficiency theory then characterized efficient influence functions, regular asymptotically linear estimators, and one-step debiased procedures \citep{levit1975efficiency,pfanzagl1985contributions,hasminskii1979nonparametric,bickel1993efficient}. In missing-data and causal-inference problems, these ideas led to influence-function-based estimating equations \citep{robinsCausal,robins1995analysis,vanderlaanunified,tsiatis2006semiparametric} and doubly robust estimation \citep{bang2005doubly}. From the perspective of asymptotic efficiency, missing-data models under missingness or coarsening at random are by now well understood: the class of regular asymptotically linear estimators can be characterized through their influence functions in these models \citep{robinsCausal,robins1995analysis,vanderlaanunified} (see Appendix \ref{appendix::momentgeneral} for a review).  

 
 
\textbf{Debiased machine learning.} A large body of work combines the semiparametric framework above with machine-learning-based nuisance estimation, often under the label of debiased machine learning. Flexible learning methods in estimating equations and AIPW-style procedures already appear in \cite{vanderlaanunified}. Targeted minimum loss estimation (TMLE) \citep{van2006targeted,vanderLaanRose2011} provides a closely related framework in which an initial nuisance estimator is updated in a parameter-targeted way to debias the resulting plug-in estimator; in the present setting, the main nuisance is the outcome regression \citep{gruber2009targeted,hoffman2020tmle,ross2025constructing}. These approaches are often combined with sample splitting and cross-fitting to accommodate generic machine learning estimators \citep{schick1986asymptotically,zheng2010asymptotic,vanderLaanRose2011,DoubleML}. Double machine learning emphasizes similar ingredients---orthogonality, cross-fitting, and influence-function expansions---for inference with modern nuisance estimators \citep{DoubleML,chernozhukov2022automatic}, and can be understood within the same semiparametric framework \citep{vanderlaan2019cvtmle,diaz2020machine,chen2026equivalence}. Closely related work also studies balancing-weight estimators \citep{hainmueller2012entropy,imai2014covariate,zubizarreta2015stable,lendle2015balancing,chattopadhyay2020balancing,hejazi2023revisiting}, which often admit equivalent regression-adjustment or augmented representations \citep{chattopadhyay2023implied,bruns2025augmented}. In missing-at-random settings, including semisupervised learning and randomized controlled trials, such estimators are often doubly robust, allowing flexible black-box methods to be used for nuisance estimation while preserving valid inference.

 


\textbf{Semisupervised inference.} Related work studies semisupervised or surrogate-assisted estimation in settings where the primary outcome is observed only on a limited subset, while auxiliary outcomes or surrogate measurements are available more broadly \citep{pepe1992inference,pepe1994auxiliary,chen2000miscellanea,chen2008improving,cheng2021robust,ji2025predictions,kallus2025role}. Two-sample semisupervised inference is a special case of data fusion, for which a substantial semiparametric efficiency literature has been developed \citep{li2023efficient,li2025data,graham2024towards,xu2025unified}. In this context, more recent papers have used the term \emph{prediction-powered inference} for the semisupervised mean-estimation setting in which a black-box predictor is combined with a small labeled sample and a large unlabeled sample \citep{ppi2023,angelopoulos2023ppi,crossppi2024,xu2025unified,song2026demystifying}. However, the resulting estimators are special cases of earlier debiased machine learning and semiparametric methods from the missing-data literature. In particular, PPI++ is an AIPW estimator obtained by empirical efficiency maximization \citep{rubin2008empirical} over a scaling parameter, whereas cross-PPI \citep{crossppi2024} is an instance of AIPW \citep{robins1995analysis} with cross-fitting, which accommodates machine-learning nuisance estimators and relaxes empirical process conditions, as in \citet{schick1986asymptotically,zheng2010asymptotic,DoubleML}. Its stability-based proof strategy is also similar in spirit to the analysis of \citet{chen2022stable}. More fundamentally, robustness to black-box prediction models is an immediate consequence of the double robustness of AIPW when the missingness probability is constant and known \citep{bang2005doubly}.

\textbf{From missing data to causal inference.} From the perspective of estimation and efficiency theory, the semisupervised setting with missing-at-random outcomes is statistically equivalent to randomized experiments, controlled trials, and survey-sampling designs once the missingness indicator is identified with the treatment or sampling indicator, so that the corresponding potential outcomes are treated as missing \citep{rubin1978bayesian,vanderlaanunified,licausal}. Accordingly, tools for covariate adjustment, causal inference, and model-assisted estimation developed in those settings can be imported directly into the semisupervised setting; see, for example, \citet{mozer2026ppi} for equivalence results relating prediction-powered inference to model-assisted survey-sampling estimators. Moreover, the two-sample setting considered here is statistically equivalent to a one-sample i.i.d.\ setting \citep{li2023efficient}. In particular, the methods we develop also apply to inference on the counterfactual mean \(\mathbb{E}[Y(1)]\) in a randomized controlled trial with binary treatment \(A_i \in \{0,1\}\), where \(\{(X_i, Y_i)\}_{i=1}^n\) are the observed covariates and outcomes for units assigned to treatment \(A_i=1\), so that \(Y_i = Y_i(1)\), and the unlabeled dataset \(\{\widetilde X_i\}_{i=1}^N\) consists of covariate information for units assigned to control \(A_i=0\) (see Appendix \ref{app:causal}).

\textbf{Calibration.} A final relevant thread concerns calibration. Classical methods such as Platt scaling \citep{platt1999probabilistic,cox1958two}, isotonic regression \citep{zadrozny2002transforming,niculescu2005predicting}, and histogram binning \citep{zadrozny2001obtaining} improve predictive reliability without retraining the underlying model. The notion of calibration appears to trace back to \citet{mincer1969evaluation}; see also \citet{lichtenstein1977calibration} and \citet{lee2023t} for historical discussion. More recent work studies calibration under broader losses and prediction tasks \citep{noarov2023statistical,van2023causal,whitehouse2024orthogonal,van2025venn}. Calibration has also been used in conformal prediction \citep{van2024self,van2025venn}, importance-weight stabilization \citep{van2024stabilized}, treatment effect estimation \citep{chernozhukov2018generic,van2023causal,whitehouse2024orthogonal}, and value prediction in offline reinforcement learning \citep{van2025bellman}. Most closely related is recent work connecting calibration to semiparametric inference by formulating it through score equations and using it to support debiased estimation of statistical functionals \citep{van2024automatic,van2023adaptive,van2025semiparametric}. In particular, \citet{van2024automatic} studies calibrated DML for doubly robust inference, which in the present setting reduces to calibrating the outcome regression and exploiting the resulting plug-in representation, while \citet{van2025semiparametric} studies calibrated plug-in regression estimators in reinforcement learning. More broadly, \citet{van2023adaptive} studies plug-in calibration for both regression and treatment effect estimation, and shows how calibration can support superefficient inference through dimension reduction.

 
 
 

\section{Two-Sample Semisupervised Setup}
 \label{sec:setup}
\subsection{Data structure and notation}
We observe two independent datasets:
\[
\mathcal D_L=\{(X_i,Y_i)\}_{i=1}^n
\qquad\text{and}\qquad
\mathcal D_U=\{\widetilde X_j\}_{j=1}^N.
\]
The labeled sample is i.i.d.\ from a distribution \(P_0\) on \(\mathcal X \times \R\), while the unlabeled sample is i.i.d.\ from the marginal distribution \(\Px\) of \(X\) under \(P_0\). We target the population mean
\[
\psi_0 := \E_{P_0}[Y],
\qquad
\muzero(x):=\E_{P_0}[Y \mid X=x].
\]

It is helpful to keep the joint two-sample experiment explicit. For each pair \((n,N)\), the data law is
\[
\mathbb D_0^{(n,N)} := P_0^{\otimes n}\otimes \Px^{\otimes N},
\]
and the target parameter is the functional
\[
\Psi\left(\mathbb D_0^{(n,N)}\right)
:=
\psi_0
=
\int y\,dP_0(x,y)
=
\int \muzero(x)\,d\Px(x).
\]
For measurable functions \(a:\mathcal X\times \R \to \R\) and \(b:\mathcal X \to \R\), define the empirical means
\[
\Pl a := \frac{1}{n}\sum_{i=1}^n a(X_i,Y_i),
\qquad
\Pu b := \frac{1}{N}\sum_{j=1}^N b(\widetilde X_j).
\]
For clarity, we will sometimes abuse notation and write, for example, \(\Pl\{a(X,Y)\}\) for \(\Pl a\) and, when a function depends only on \(X\), write \(\Pl \{b(X)\} := \frac{1}{n}\sum_{i=1}^n b(X_i)\).
Throughout, we write \(\rho_n = n/M\) for the labeled fraction, equivalently the known probability of observing the outcome, and use the identity
\[
\frac{1}{n+N}\left\{\sum_{i=1}^n f(X_i) + \sum_{j=1}^N f(\widetilde X_j)\right\}
=
\rho_n \Pl f + (1-\rho_n)\Pu f.
\]
That is, the empirical mean of \(f(X)\) over the pooled covariate sample is a convex combination of the empirical means in the labeled and unlabeled samples.

 


\begin{assumption}[Two-sample design]
\label{assump:design}
The following hold:
\begin{enumerate}[leftmargin=6mm,label=(\roman*)]
\item \(\mathcal D_L\) is i.i.d.\ from \(P_0\), and \(\mathcal D_U\) is i.i.d.\ from \(\Px\).
\item \(\mathcal D_L\) and \(\mathcal D_U\) are independent.
\item \(Y\) has finite second moment under \(P_0\).
\item With \(M:=n+N\), the labeled fraction \(\rhon:=n/M\) satisfies \(\rhon \to \rhozero\) for some \(\rhozero \in (0,1)\).
\end{enumerate}
\end{assumption}

\begin{remark}
The labeled sample carries the outcome information, while the unlabeled sample sharpens estimation of the marginal covariate distribution. The pooled i.i.d.\ missing-data formulation in the appendix is a special case of this two-sample setup. 
\end{remark}


\subsection{Review of AIPW, PPI, and semiparametric efficiency}
\label{sec:effreview}

For later reference, we briefly review semiparametric efficiency in the unrestricted two-sample model \citep{bickel1993efficient,li2023efficient}, along with the associated class of influence functions and AIPW estimators \citep{robinsCausal,robins1995analysis,bang2005doubly}.   We say that an estimator \(\widehat\psi\) is asymptotically linear with \emph{influence pair} \((D^L,D^U)\) if
\[
\sqrt{M}\,(\widehat\psi-\psi_0)
=
\frac{1}{\sqrt{M}}\sum_{i=1}^n D^L(X_i,Y_i)
+
\frac{1}{\sqrt{M}}\sum_{j=1}^N D^U(\widetilde X_j)
+
o_p(1),
\]
where \(\mathbb{E}_{P_0}[D^L(X,Y)]=0\) and \(\mathbb{E}_{\Px}[D^U(\widetilde X)]=0\). In other words, the error \(\widehat\psi-\psi_0\) is asymptotically equivalent to the empirical mean \(\frac{1}{M}\{\sum_{i=1}^n D^L(X_i,Y_i)+\sum_{j=1}^N D^U(\widetilde X_j)\}\). This is the natural two-sample analogue of the usual influence function in an i.i.d.\ model \citep{van2000asymptotic} and uniquely determines the corresponding influence function in the i.n.i.d.\ model \citep{bickel1993efficient}. By the central limit theorem and the independence of the two datasets,
\[
\sqrt{M}\,(\widehat\psi-\psi_0)
\;\rightsquigarrow\;
N\!\left(0,\,
\rhozero \Var_{P_0}\bigl(D^L(X,Y)\bigr)
+
(1-\rhozero)\Var_{\Px}\bigl(D^U(\widetilde X)\bigr)\right).
\]
Thus, the influence pair \((D^L,D^U)\) determines the asymptotic variance of \(\widehat\psi\).

The following result characterizes the class of regular estimators \citep{van2000asymptotic} in this model through their influence functions and shows that each is asymptotically equivalent to an AIPW estimator. This provides the benchmark class for our later analysis, where we place calibrated plug-in estimators within the standard AIPW framework and characterize their efficiency. The result follows from \citet{robins1995analysis} together with the statistical equivalence between the two-sample model and the single-sample i.i.d.\ model established, for example, by \citet{li2023efficient}.

\begin{proposition}[AIPW class of influence pairs]
\label{prop:if-class}
Under \cref{assump:design}, a mean-zero pair \((D^L,D^U)\) is the influence pair of a regular estimator of \(\psi_0\) in the unrestricted two-sample model if and only if there exists a square-integrable function \(f:\mathcal X\to\R\) such that, with
\[
\tilde f(X):=f(X)-\E_{\Px}[f(X)]+\psi_0,
\]
the pair takes the form
\begin{align}
\label{eq:if-class-L}
D_f^L(X,Y)
&=
\tilde f(X)-\psi_0+\rhozero^{-1}\{Y-\tilde f(X)\},\\
\label{eq:if-class-U}
D_f^U(\widetilde X)
&=
\tilde f(\widetilde X)-\psi_0.
\end{align}
Furthermore, for any such \(f:\mathcal X\to\R\), a corresponding estimator with this influence pair is given by the AIPW estimator
\begin{equation}
\label{eq:aipw-family}
\widehat\psi(f)
:=
\rhon \Pl\{f(X)\}
+
(1-\rhon)\Pu\{f(\widetilde X)\}
+
\Pl\{Y-f(X)\}.
\end{equation}
\end{proposition}
 

\textbf{AIPW, PPI, and PPI++.}
A classical choice is the AIPW estimator \citep{robins1995analysis} obtained from the plug-in function \(f(X)=m(X)\), where \(m(X)\) denotes the prediction score, namely
\[
\widehat{\psi}_{\mathrm{AIPW}}
=
\rhon \Pl\{m(X)\}
+
(1-\rhon)\Pu\{m(\widetilde X)\}
+
\Pl\{Y-m(X)\}.
\]
This can be written in the usual \textit{iid} form by defining the pooled empirical measure \(\mathbb{P}_{M}:=\rhon \Pl + (1-\rhon)\Pu\), in which case
\[
\widehat{\psi}_{\mathrm{AIPW}}
=
\mathbb{P}_{M}\{m(X)\}
+
\rhon^{-1}\mathbb{P}_{M}\{Y-m(X)\},
\]
where \(\rhon^{-1}\) is the known inverse probability of observing the outcome. Since outcome missingness is completely at random, AIPW estimators are ``safe'' in the sense that inference remains valid regardless of the quality of \(m\) \citep{xu2025unified}; \(m\) affects only the variance of the estimator. This is a special case of double robustness: because the missingness probability \(\rho_n\) is constant and known, the estimator remains unbiased even if \(m\) is misspecified \citep{bang2005doubly}.

An alternative choice takes \(f(X)=m(X)/(1-\rho_n)\), which yields
\[
\widehat{\psi}_{\mathrm{PPI}}
=
\Pu\{m(\widetilde X)\}+\Pl\{Y-m(X)\}.
\]
This choice is used implicitly in \citet{ppi2023}, where it is termed PPI\footnote{The PPI and PPI++ estimators are algebraically equivalent to the model-assisted survey-design estimators of \citet{cassel1976some} and \citet{sarndal2003model}, respectively; see \citet{mozer2026ppi} for a review. In Section~\ref{sec:linear}, we further show that PPI++ is equivalent to linear regression adjustment using \(m(X)\) as a prognostic score.}. Relative to \(\widehat{\psi}_{\mathrm{AIPW}}\), \(\widehat{\psi}_{\mathrm{PPI}}\) omits the labeled-sample average of \(m(X)\) from the plug-in term. It is therefore typically less efficient, especially when \(m\) is reasonably accurate; in particular, it is inefficient when \(m=\mu_0\). In the regime where \(\rho_n \approx 0\) and \(m\) is inaccurate, PPI can be more efficient when \(\Cov(m(X), Y-m(X))\) is sufficiently positive, since it avoids the cross-correlation between the plug-in term and the bias-correction term.

 

Because both \(\widehat{\psi}_{\mathrm{AIPW}}\) and \(\widehat{\psi}_{\mathrm{PPI}}\) correspond to choices in the scaling class \(\{\lambda m(X):\lambda\in\mathbb R\}\), one can construct an estimator that is asymptotically at least as efficient as both by minimizing the empirical influence-function variance criterion
\[
\sigma_n^2(f)
:=
\rho_n \Pl\left[\left\{f(X)-\widehat{\psi}(f)+\rho_n^{-1}(Y-f(X))\right\}^2\right]
+
(1-\rho_n)\Pu\left[\left\{f(\widetilde X)-\widehat{\psi}(f)\right\}^2\right]
\]
over \(\{\lambda m(X): \lambda \in \mathbb{R}\}\), choosing \(\lambda\) by empirical efficiency maximization \citep{vanderlaanunified,rubin2008empirical}. The resulting estimator, termed PPI++ in \citet{angelopoulos2023ppi}, is simply AIPW with an optimally chosen multiple of the prediction score. Efficiency can be improved further by optimizing over richer classes, such as
\[
\{\lambda_{\mathrm{shift}} + \lambda_{\mathrm{scale}} m(X) + \beta_{\mathrm{lin}}^\top X : \lambda_{\mathrm{shift}}, \lambda_{\mathrm{scale}} \in \mathbb{R},\ \beta_{\mathrm{lin}} \in \mathbb{R}^d\}.
\]


The semiparametrically efficient benchmark is attained by the unrestricted choice \(f(X)=\mu_0(X)\), which must typically be estimated from the data. From the balancing perspective, this shows that efficiency generally requires using the full covariate information \(X\), or at least a summary rich enough to recover the true regression function \(\mu_0(X)\) \citep{robins1995analysis}. The next result identifies the efficient influence pair, and thus the most efficient regular estimator, in this model \citep{bickel1993efficient}. Accordingly, efficiency is optimized by using a prediction model \(m(X)\) that approximates \(\mu_0(X)\) as closely as possible. Modern debiased machine learning methods implement this principle using flexible black-box learners to estimate nuisance functions \citep{vanderlaanunified,vanderLaanRose2011,DoubleML}.



\begin{proposition}[Efficient influence pair in the full model]
\label{prop:eif-full}
Under \cref{assump:design}, the asymptotic variance over the class \cref{eq:if-class-L,eq:if-class-U} is minimized at \(f=\muzero\). The efficient influence pair for \(\psi_0=\E_{P_0}[Y]\) in the unrestricted two-sample model is therefore
\begin{align}
\label{eq:eif-full-L}
D_{\mathrm{eff}}^L(X,Y)
&=
\muzero(X)-\psi_0+\rhozero^{-1}\{Y-\muzero(X)\},\\
\label{eq:eif-full-U}
D_{\mathrm{eff}}^U(\widetilde X)
&=
\muzero(\widetilde X)-\psi_0.
\end{align}
\end{proposition}
 
  

To conclude, we note that regression adjustment, balancing-weight adjustment, and AIPW estimation are closely related, as the following remark illustrates.

\begin{remark}[Plug-in approaches to bias correction and TMLE]
\label{remark::plugin}
The same first-order class can also be realized through regression-based or balancing-weight adjustment \citep{scharfstein1999adjusting,van2006targeted,rosenblum2009using}. Let
\[
m_n(X)=\hat a_{\mathrm{lin}}+\hat b_{\mathrm{lin}}\,m(X),
\]
where
\[
(\hat a_{\mathrm{lin}},\hat b_{\mathrm{lin}})
=
\arg\min_{a,b}\sum_{i=1}^n \{Y_i-a-b\,m(X_i)\}^2.
\]
Then, for any \(\lambda\in\mathbb{R}\),
\[
\rho_n \Pl\{m_n(X)\}
+
(1-\rho_n)\Pu\{m_n(\widetilde X)\}
=
\rho_n \Pl\{m_n(X)\}
+
(1-\rho_n)\Pu\{m_n(\widetilde X)\}
+
\lambda \Pl\bigl[Y-m_n(X)\bigr],
\]
because the augmentation term vanishes by the normal equations for the least-squares fit. Dually, consider affine weights \(w_n(X)=\hat a_{\mathrm{bal}}+\hat b_{\mathrm{bal}}\,m(X)\) that balance on \((1,m(X))\), in the sense that \(\Pl\{w_n(X)\}=1\) and \(\Pl\bigl[m(X)w_n(X)\bigr]=\rho_n \Pl\{m(X)\}+(1-\rho_n)\Pu\{m(\widetilde X)\}\). The corresponding weighted estimator satisfies
\[
\Pl\bigl\{w_n(X)Y\bigr\}
=
\rho_n \Pl\{m(X)\}
+
(1-\rho_n)\Pu\{m(\widetilde X)\}
+
\Pl\bigl[w_n(X)\{Y-m(X)\}\bigr].
\]
In fact,
\[
\Pl\bigl\{w_n(X)Y\bigr\}
=
\rho_n \Pl\{m_n(X)\}
+
(1-\rho_n)\Pu\{m_n(\widetilde X)\}.
\]
Thus, the plug-in estimator based on \(m_n\), the balancing-weight estimator based on \(w_n\), and the AIPW estimator based on \((\lambda,m_n)\) are all algebraically equivalent; see \citet{chattopadhyay2023implied,bruns2025augmented} for related equivalence results. Consequently, whenever \(m_n\) converges to the best linear predictor of \(Y\) given \(m(X)\), these estimators are typically at least as efficient as both \(\widehat{\psi}_{\mathrm{AIPW}}\) and \(\widehat{\psi}_{\mathrm{PPI}}\), since they correspond to an influence function with smaller or equal variance. In fact, they are algebraically equivalent to the PPI++ estimator; see Appendix~\ref{sec::ppi++} and Section \ref{sec:linear}.

More broadly, such debiased plug-in estimators can be viewed as instances of TMLE \citep{van2006targeted,vanderLaanRose2011}, in which either the weights or the regression model are updated so that the corresponding augmentation term vanishes. In weighting-based TMLE, the update balances on the prediction model \(m(X)\) (e.g., \citealt{hejazi2023revisiting}); one may either include the known weight \(1/\rho_n\) as an offset or balance on the intercept as well. In regression-based TMLE, by contrast, the prediction model \(m(X)\) typically enters as an offset in the linear model rather than as a covariate (e.g., \citealt{lendle2015balancing}); the intercept adjustment then plays the role of correcting for the known inverse missingness-probability weight \(1/\rho_n\).
\end{remark}





\section{Calibrated Prediction-Powered Inference}
\label{sec:calppi}
The first subsection introduces the general calibrated plug-in estimator and shows that it admits an equivalent AIPW representation \citep{chattopadhyay2023implied, bruns2025augmented}. The second subsection then describes concrete ways to enforce empirical calibration, focusing on isotonic regression. We next discuss linear calibration as a simpler adjustment step when data are especially scarce. The final subsection presents a cross-fitted version of the same construction, which allows the predictive score to be learned from the labeled data.

\subsection{General estimator and exact AIPW representation}

Let \(\mhat:\mathcal X\to\mathbb R\) denote a black-box regression fit on the labeled data. A calibration procedure takes the labeled scores \(\mhat(X_1),\dots,\mhat(X_n)\) and outcomes \(Y_1,\dots,Y_n\), fits a post-hoc map \(\fhat\), and outputs the calibrated regression
\[
\mtilde := \fhat\circ\mhat.
\]
Thus, the fitted values on the two samples are \(\mtilde(X_i)\) for \(i=1,\dots,n\) and \(\mtilde(\widetilde X_j)\) for \(j=1,\dots,N\).

Fix a linear class \(\mathcal F\) of real-valued functions \(f:\mathbb R\to\mathbb R\) that contains the constant functions $t \mapsto c$ and the identity map \(t\mapsto t\). Informally, calibration is a reliability property stating that the prediction model \(\mtilde\) cannot be improved by post-processing its output using transformations in \(\mathcal F\). We say that the labeled calibrated predictions are \emph{empirically calibrated over \(\mathcal F\)} if no admissible transformation of the calibrated score improves the empirical squared risk, that is,
\begin{equation}
\label{eq:cal-risk}
\Pl\bigl[(Y-\mtilde(X))^2\bigr]
\le
\Pl\bigl[(Y-f(\mtilde(X)))^2\bigr]
\end{equation}
for every \(f\in\mathcal F\). Equivalently, the corresponding first-order optimality conditions imply
\begin{equation}
\label{eq:cal-score}
\Pl\bigl[f(\mtilde(X))\{Y-\mtilde(X)\}\bigr]=0,
\qquad
f\in\mathcal F.
\end{equation}
Because \(\mathcal F\) contains the constant functions, empirical calibration in particular implies the mean-calibration identity
\[
\Pl\{Y-\mtilde(X)\}=0.
\]

The empirically calibrated plug-in estimator is
\begin{equation}
\label{eq:plugin}
\psihat
:=
\frac{1}{n+N}\left\{\sum_{i=1}^n \mtilde(X_i) + \sum_{j=1}^N \mtilde(\widetilde X_j)\right\}
=
\rhon \Pl\{\mtilde(X)\}+(1-\rhon)\Pu\{\mtilde(\widetilde X)\},
\end{equation}
where \(\Pl\) and \(\Pu\) denote the empirical means over the labeled and unlabeled samples. 


We focus on two cases. In the nonparametric case, \(\mathcal F\) is the class of all transformations \(f:\R\to\R\), which we refer to as \emph{full calibration}; in practice, this can be implemented by regressing \(Y\) on \(m(X)\) using histogram regression or isotonic regression (Section \ref{sec:iso}). In the parametric case, \(\mathcal F\) is restricted to affine transformations \(t\mapsto at+b\), which we refer to as \emph{linear calibration}; this amounts to fitting \((a,b)\) in the working linear model \(a\,m(X)+b\) for \(Y\) (Section \ref{sec:linear}). More generally, if \(\mathcal F\) is finite-dimensional with basis \(\{\varphi_1,\dots,\varphi_p\}\), then an empirically calibrated predictor \(m_n^*\) can be constructed by regressing \(Y\) on \(\{\varphi_1(m(X)),\dots,\varphi_p(m(X))\}\)\footnote{If \(\mathcal F\) does not contain the identity map, one can instead include \(m(X)\) as an offset, and \eqref{eq:cal-score} will still hold.}.

The key property of calibrated plug-in estimators is the following:
 

\begin{theorem}[Exact AIPW representation under empirical calibration]
\label{thm:repr}
Suppose the calibrated predictions satisfy \cref{eq:cal-score}, where \(\mathcal F\) is a linear class containing the constant functions and the identity map. Then, for any empirical balancing weight \(\what\in\mathcal F\), the calibrated plug-in estimator admits the exact AIPW representation
\begin{equation}
\label{eq:abw-general}
\psihat
=
\rhon \Pl\{\mtilde(X)\}+(1-\rhon)\Pu\{\mtilde(\widetilde X)\}
+
\rhon \Pl\bigl[\what(\mtilde(X))\{Y-\mtilde(X)\}\bigr].
\end{equation}
Moreover, if \(\what_n^*\in\mathcal F\) balances \(\mathcal F\) in the sense that
\begin{equation}
\label{eq:F-balance}
\Pl\bigl[\what_n^*(\mtilde(X))\,f(\mtilde(X))\bigr]
=
\rhon \Pl\{f(\mtilde(X))\}
+
(1-\rhon)\Pu\{f(\mtilde(\widetilde X))\},
\qquad
f\in\mathcal F,
\end{equation}
then \(\psihat\) admits the weighted-outcome representation
\begin{equation}
\label{eq:abw-general-balance}
\psihat
=
\Pl\bigl[\what_n^*(\mtilde(X))Y\bigr].
\end{equation}
\end{theorem}

Theorem \ref{thm:repr} shows that the calibrated plug-in estimator \(\psihat\) is identical to the AIPW estimator with \(\mtilde\) as the prediction score. It also yields a dual representation: \(\psihat\) can be written as a balancing-weighted outcome estimator, where \(\what_n^*\) is chosen to balance moments of functions in \(\mathcal F\) between the labeled and unlabeled samples \citep{chattopadhyay2020balancing}.

\textbf{Sketch of asymptotic theory.} We now sketch the asymptotic theory for \(\psihat\). Suppose the estimated calibrated score \(\mtilde\) converges to a limit \(m_0\), typically a transformation of \(m\) that minimizes a population risk over a suitable class. This limit is typically calibrated at the population level in the sense that
\[
\mathbb{E}_{P_0}\bigl[f\{m_0(X)\}\{\mu_0(X)-m_0(X)\}\bigr]=0,
\qquad
f\in\mathcal F.
\]
If \(\mathcal F\) contains all real-valued functions, then
\[
m_0(X)=\mathbb{E}_{P_0}\left[Y\mid m_0(X)\right],
\]
so the predictive performance of \(m_0\) cannot be improved by any transformation of its output. Taking \(\hat w=1\) in \eqref{eq:abw-general}, \(\psihat\) can be viewed as a plug-in version of the usual AIPW estimator with \(\mtilde\) as the regression adjustment. Therefore, if \(\mtilde\) is sufficiently close to \(m_0\), then \(\psihat\) should behave asymptotically like the corresponding oracle AIPW estimator:
\begin{align*}
\psihat-\psi_0
&=
\rhon(\Pl-P_0)\Bigl\{m_0(X)-\psi_0+\rhon^{-1}\bigl(Y-m_0(X)\bigr)\Bigr\}\\
&\quad+(1-\rhon)(\Pu-\Px)\bigl\{m_0(\widetilde X)-\psi_0\bigr\}
+o_p(M^{-1/2}).
\end{align*}
In particular, \(\psihat\) is asymptotically linear with influence pair
\[
D_{m_0}^L(X,Y)=m_0(X)-\psi_0+\rhozero^{-1}\{Y-m_0(X)\},
\qquad
D_{m_0}^U(\widetilde X)=m_0(\widetilde X)-\psi_0.
\]
In Section~\ref{sec:theory}, we establish this result for the isotonic-calibrated plug-in estimator introduced next.

\subsection{Achieving full calibration via isotonic regression}
\label{sec:iso}

In practice, full empirical calibration can be achieved by fitting a second-stage calibrator to the labeled pairs \(\{(\mhat(X_i),Y_i)\}_{i=1}^n\) \citep{van2025venn}. Two especially simple options are histogram binning and isotonic calibration \citep{zadrozny2001obtaining,zadrozny2002transforming,niculescu2005predicting}, both of which yield exact empirical calibration. Histogram binning fits a piecewise-constant map on a prespecified partition of the score \citep{gupta2021distribution}, whereas isotonic calibration estimates the best monotone transformation of the score and adaptively determines the effective bins using the pooled-adjacent-violators algorithm \citep{barlow1972isotonic}. We focus on isotonic regression because it is lightweight, tuning-free, and easy to implement using widely available software. Moreover, because the identity map is monotone, isotonic calibration can only improve the empirical mean squared error over \(\mathcal{D}_L\) relative to the original scores.


The isotonic-calibrated plug-in estimator is
\[
\widehat{\psi}_{\mathrm{iso}}
:=
\frac{1}{n+N}
\left\{
\sum_{i=1}^n m_{n,\mathrm{iso}}^\star(X_i)
+
\sum_{j=1}^N m_{n,\mathrm{iso}}^\star(\widetilde X_j)
\right\},
\]
where
\[
m_{n,\mathrm{iso}}^\star(X)
=
\hat f\{\mhat(X)\},
\qquad
\hat f \in \arg\min_{f \in \mathcal F_{\mathrm{iso}}}
\sum_{i=1}^n \bigl\{Y_i - f(\mhat(X_i))\bigr\}^2,
\]
and \(\mathcal F_{\mathrm{iso}}\) denotes the class of nondecreasing real-valued functions. The fitted values \(m_{n,\mathrm{iso}}^\star(X_i)=\hat f(\mhat(X_i))\) are blockwise empirical means determined by the pooled-adjacent-violators algorithm. This blockwise structure implies full empirical calibration, even though the recalibration step is restricted to isotonic functions.  

Algorithm~\ref{alg:isoppi} summarizes the resulting estimator in the same ``fit, debias, infer'' style commonly used in the PPI literature.

\begin{algorithm}[t]
\caption{Isotonic-calibrated plug-in estimator}
\label{alg:isoppi}
\begin{algorithmic}[1]
\Require \(\mathcal D_L=\{(X_i,Y_i)\}_{i=1}^n\), \(\mathcal D_U=\{\widetilde X_j\}_{j=1}^N\), score \(\mhat\), level \(1-\alpha\)
\State Compute \(T_i=\mhat(X_i)\) and \(\widetilde T_j=\mhat(\widetilde X_j)\)
\State Fit \(\hat f\in\arg\min_{f\in\mathcal F_{\mathrm{iso}}}\sum_{i=1}^n\{Y_i-f(T_i)\}^2\)
\State Set \(m_{n,\mathrm{iso}}^\star(X_i)=\hat f(T_i)\) and \(m_{n,\mathrm{iso}}^\star(\widetilde X_j)=\hat f(\widetilde T_j)\)
\State Form \(\widehat\psi_{\mathrm{iso}}\gets (n+N)^{-1}\{\sum_{i=1}^n m_{n,\mathrm{iso}}^\star(X_i)+\sum_{j=1}^N m_{n,\mathrm{iso}}^\star(\widetilde X_j)\}\)
\State Set \(\widehat D_i^L=m_{n,\mathrm{iso}}^\star(X_i)-\widehat\psi_{\mathrm{iso}}+\rhon^{-1}\{Y_i-m_{n,\mathrm{iso}}^\star(X_i)\}\) and \(\widehat D_j^U=m_{n,\mathrm{iso}}^\star(\widetilde X_j)-\widehat\psi_{\mathrm{iso}}\)
\State Set \(\widehat{\mathrm{SE}}_{\mathrm{iso}}^2=(n+N)^{-2}\{\sum_{i=1}^n(\widehat D_i^L)^2+\sum_{j=1}^N(\widehat D_j^U)^2\}\)
\State Return \(\widehat\psi_{\mathrm{iso}} \pm z_{1-\alpha/2}\widehat{\mathrm{SE}}_{\mathrm{iso}}\)
\end{algorithmic}
\end{algorithm}
 

\begin{proposition}[Isotonic calibration is fully calibrated]
\label{prop:iso-repr}
For every function \(h:\R\to\R\),
\begin{equation}
\label{eq:orth}
\Pl\bigl[h(m_{n,\mathrm{iso}}^\star(X))\{Y-m_{n,\mathrm{iso}}^\star(X)\}\bigr]=0.
\end{equation}
Moreover, it holds that
\[
\Pl\bigl[(Y-m_{n,\mathrm{iso}}^\star(X))^2\bigr]
\le
\Pl\bigl[(Y-\mhat(X))^2\bigr].
\]
\end{proposition}

Consequently, the AIPW representation of Theorem \ref{thm:repr} holds with \(\mathcal F\) equal to the class of all real-valued functions. In the context of treatment effect calibration, Proposition \ref{prop:iso-repr} was derived by \citet{van2023causal}.

 
 


\begin{remark}[Binning calibrators achieve full calibration]
Histogram binning and related piecewise-constant calibrators, which assign recalibrated predictions as the empirical mean of the labels within each bin, also satisfy Proposition~\ref{prop:iso-repr} \citep{van2025venn}. We do not treat them separately, since the isotonic case already captures the main argument.
\end{remark}

\begin{remark}[Implementation]
Isotonic regression is widely available in standard software. A standard algorithm is the pooled-adjacent-violators algorithm (PAVA) \citep{barlow1972isotonic}. Although it is generally fast, it may become slow at ultra-large sample sizes. A scalable alternative is to use regression tree software with monotonicity constraints, such as XGBoost or LightGBM (see Appendix \ref{appendix:code}). In one dimension, isotonic regression is equivalent to fitting a regression tree of infinite depth under a monotonicity constraint. An advantage of tree-based software is that the fit can be regularized by constraining or tuning the number of leaves, the maximum tree depth, or the minimum number of observations per leaf. 
\end{remark}


\begin{remark}[Shrinkage via Venn--Abers]
When the labeled sample is small, isotonic calibration may be unstable. A natural alternative is Venn--Abers calibration, originally introduced for binary classification \citep{vovk2012venn}; see \cite{van2024self, van2025venn} for generalized versions beyond the binary setting. Assume \(Y \in [0,1]\), or rescale otherwise. For each \(x\), augment the labeled calibration sample with the hypothetical point \((\mhat(x),y)\) for \(y \in \{0,1\}\), fit the resulting isotonic calibrators \(f_n^{(x,0)}\) and \(f_n^{(x,1)}\), and form the interval-valued prediction
\[
\bigl[f_n^{(x,0)}\{\mhat(x)\},\, f_n^{(x,1)}\{\mhat(x)\}\bigr].
\]
Following \cite{vovk2012venn}, one may convert this interval to a point prediction by shrinking its midpoint toward the (unadjusted) AIPW estimator:
\[
m_{n,\mathrm{VA}}^*(x)
:=
m_{n,\mathrm{mid}}^*(x)
+
\bigl\{f_n^{(x,1)}\{\mhat(x)\}-f_n^{(x,0)}\{\mhat(x)\}\bigr\}
\bigl\{\widehat\psi_{\mathrm{unadj},0}-m_{n,\mathrm{mid}}^*(x)\bigr\},
\]
where
\[
\begin{aligned}
m_{n,\mathrm{mid}}^*(x)
&:=
\frac{f_n^{(x,1)}\{\mhat(x)\}+f_n^{(x,0)}\{\mhat(x)\}}{2},\\
\widehat\psi_{\mathrm{unadj},0}
&:=
\frac{1}{n+N}
\left\{
\sum_{i=1}^n \mhat(X_i)
+
\sum_{j=1}^N \mhat(\widetilde X_j)
\right\}
+
\frac{1}{n}\sum_{i=1}^n \{Y_i-\mhat(X_i)\}.
\end{aligned}
\]
Since the shrinkage magnitude is proportional to the width of the Venn--Abers interval, this construction can be more stable than plain isotonic calibration in small samples.
\end{remark}
 
\subsection{Achieving linear calibration via linear regression}
\label{sec:linear}
For very small labeled samples, one can replace histogram binning and isotonic calibration with a smooth or affine correction of the score \citep{jiang2011smooth,lucena2018spline}. Here we focus on linear calibration. This imposes a more restrictive calibration model than isotonic regression, similar to Platt scaling and temperature scaling \citep{platt1999probabilistic,guo2017calibration}, but can be more stable in small samples. Let
\begin{equation}
\label{eq:linear-cal}
\bigl(\hat a_{\mathrm{lin}},\hat b_{\mathrm{lin}}\bigr)\in\arg\min_{a,b\in\R}\sum_{i=1}^n \{Y_i-a\,\mhat(X_i)-b\}^2,
\qquad
m_{n,\mathrm{lin}}^\star(X):=\hat a_{\mathrm{lin}}\,\mhat(X)+\hat b_{\mathrm{lin}}.
\end{equation}
Define the associated plug-in estimator\footnote{We already saw this estimator in Remark~\ref{remark::plugin}.}by
\begin{equation}
\label{eq:linear-plugin}
\widehat{\psi}_{\mathrm{lin}}
:= \frac{1}{n+N}
\left\{
\sum_{i=1}^n m_{n,\mathrm{lin}}^\star(X_i)
+
\sum_{j=1}^N m_{n,\mathrm{lin}}^\star(\widetilde X_j)
\right\}.
\end{equation}
Because the affine class contains the constant function, this already enforces empirical mean calibration.  
 


\begin{proposition}[Linear calibration as AIPW]
\label{prop:linear}
The normal equations for \cref{eq:linear-cal} imply
\[
\Pl\{Y-m_{n,\mathrm{lin}}^\star(X)\}=0
\qquad\text{and}\qquad
\Pl\bigl[m_{n,\mathrm{lin}}^\star(X)\{Y-m_{n,\mathrm{lin}}^\star(X)\}\bigr]=0.
\]
Hence, for any \(\hat w(X)=a+b\,m_{n,\mathrm{lin}}^\star(X)\),
\begin{equation}
\label{eq:linear-ppi}
\widehat{\psi}_{\mathrm{lin}}
=
\rhon \Pl\{m_{n,\mathrm{lin}}^\star(X)\}
+(1-\rhon)\Pu\{m_{n,\mathrm{lin}}^\star(\widetilde X)\}
+
\Pl\bigl[\hat w(X)\{Y-m_{n,\mathrm{lin}}^\star(X)\}\bigr].
\end{equation}
Thus, the linearly calibrated plug-in estimator admits the usual AIPW residual-correction form.
\end{proposition}


\begin{algorithm}[t]
\caption{Linearly calibrated plug-in estimator}
\label{alg:linearppi}
\begin{algorithmic}[1]
\Require \(\mathcal D_L=\{(X_i,Y_i)\}_{i=1}^n\), \(\mathcal D_U=\{\widetilde X_j\}_{j=1}^N\), score \(\mhat\), level \(1-\alpha\)
\State Compute \(T_i=\mhat(X_i)\) and \(\widetilde T_j=\mhat(\widetilde X_j)\)
\State Fit \((\hat a,\hat b)\in\arg\min_{a,b}\sum_{i=1}^n\{Y_i-a-bT_i\}^2\)
\State Set \(m_{n,\mathrm{lin}}^\star(X_i)=\hat a+\hat bT_i\) and \(m_{n,\mathrm{lin}}^\star(\widetilde X_j)=\hat a+\hat b\widetilde T_j\)
\State Form \(\widehat\psi_{\mathrm{lin}}\gets (n+N)^{-1}\{\sum_{i=1}^n m_{n,\mathrm{lin}}^\star(X_i)+\sum_{j=1}^N m_{n,\mathrm{lin}}^\star(\widetilde X_j)\}\)
\State Set \(\widehat D_i^L=m_{n,\mathrm{lin}}^\star(X_i)-\widehat\psi_{\mathrm{lin}}+\rhon^{-1}\{Y_i-m_{n,\mathrm{lin}}^\star(X_i)\}\) and \(\widehat D_j^U=m_{n,\mathrm{lin}}^\star(\widetilde X_j)-\widehat\psi_{\mathrm{lin}}\)
\State Set \(\widehat{\mathrm{SE}}_{\mathrm{lin}}^2=(n+N)^{-2}\{\sum_{i=1}^n(\widehat D_i^L)^2+\sum_{j=1}^N(\widehat D_j^U)^2\}\)
\State Return \(\widehat\psi_{\mathrm{lin}} \pm z_{1-\alpha/2}\widehat{\mathrm{SE}}_{\mathrm{lin}}\)
\end{algorithmic}
\end{algorithm}

\begin{remark}[Platt scaling adjustment]
Linear recalibration can also be performed using losses other than squared error, such as logistic loss, Poisson loss, or, more generally, losses from canonical-link generalized linear models \citep{rosenblum2009using}. In particular, when \(Y\in\{0,1\}\) and \(\mhat(X)\in(0,1)\) is an initial predictor, one may apply Platt scaling \citep{cox1958two,platt1999probabilistic} by fitting a logistic regression of \(Y\) on the logit of \(\mhat(X)\):
\begin{equation}
\label{eq:platt-cal}
\bigl(\hat a_{\mathrm{platt}},\hat b_{\mathrm{platt}}\bigr)
\in
\arg\min_{a,b\in\R}
\sum_{i=1}^n
\Bigl[
-Y_i \log \sigma\bigl(a\,\mathrm{logit}(\mhat(X_i))+b\bigr)
-(1-Y_i)\log\Bigl(1-\sigma\bigl(a\,\mathrm{logit}(\mhat(X_i))+b\bigr)\Bigr)
\Bigr],
\end{equation}
with recalibrated predictor
\[
m_{n,\mathrm{platt}}^\star(X)
:=
\sigma\bigl(\hat a_{\mathrm{platt}}\,\mathrm{logit}(\mhat(X))+\hat b_{\mathrm{platt}}\bigr),
\]
where \(\sigma(t):=(1+e^{-t})^{-1}\). Platt scaling may be preferred over linear regression calibration when predictors exhibit over- or under-confidence, in which case calibration curves are often approximately sigmoidal. More generally, outcomes can be rescaled to \([0,1]\) before applying this procedure.
\end{remark}

 
\begin{remark}[PPI++ as prognostic-score adjustment by linear regression]
As discussed in Section~\ref{sec:effreview}, both PPI and PPI++ are special cases of AIPW estimators. They can also be understood as linear regression adjustments based on prognostic scores, a classical approach to improving precision through flexible covariate adjustment in small-sample settings \citep{hansen2008prognostic,rosenblum2009using,lin2013agnostic,schuler2022increasing,hojbjerre2025powering}. The linearly calibrated plug-in estimator \(\widehat{\psi}_{\mathrm{lin}}\) is algebraically equivalent to the PPI++ estimator of \citet{angelopoulos2023ppi}, obtained via empirical efficiency maximization \citep{rubin2008empirical}; see Appendix \ref{sec::ppi++}. Likewise, the PPI estimator \(\widehat{\psi}_{\mathrm{PPI}}\) studied in \citet{angelopoulos2023ppi} is algebraically equivalent to the unlabeled-data plug-in estimator \(\mathbb{P}_N^U\{m_{n,\mathrm{const}}^\star(X)\}\) based on the intercept-only adjustment
\begin{equation}
\label{eq:const-cal}
\hat a_{\mathrm{const}} \in \arg\min_{a\in\R}\sum_{i=1}^n \{Y_i-a-\mhat(X_i)\}^2,
\qquad
m_{n,\mathrm{const}}^\star(X):=\hat a_{\mathrm{const}}+\mhat(X),
\end{equation}
which calibrates only the mean. By contrast, the standard AIPW estimator \(\widehat{\psi}_{\mathrm{AIPW}}\) corresponds to
\[
\rho_n \mathbb{P}_n^L\{m_{n,\mathrm{const}}^\star(X)\} + (1-\rho_n)\mathbb{P}_N^U\{m_{n,\mathrm{const}}^\star(X)\}.
\]
\end{remark}


\begin{remark}[Connection with TMLE and balancing]
More broadly, our estimator can be viewed as a simple instance of TMLE \citep{van2006targeted,lendle2015balancing,hejazi2023revisiting,van2025automatic}, and is also closely connected to balancing-weight and regression-adjustment approaches, as we discuss in Remark \ref{remark::plugin}. Related equivalence results between linear regression, AIPW, and balancing-weight estimators are given by \citet{chattopadhyay2023implied,bruns2025augmented}; see also \citet{lee2026mec} for a related balancing-based estimator.
\end{remark}

 


\subsection{Cross-Fitting and Cross-Calibration}
\label{sec:crossfit}

In the main exposition, we assume for simplicity that the initial black-box predictor is fixed. In practice, however, it may be trained on the labeled sample itself, in which case the same labeled observations are reused both to fit the predictor and to calibrate its scores. A cross-fitted version avoids this reuse by replacing the single prediction function with out-of-fold predictions and then fitting the calibrator on the resulting labeled out-of-fold scores.

\begin{enumerate}[leftmargin=6mm,label=\arabic*.]
\item Partition the labeled sample indices into folds \(I_1,\dots,I_K\).
\item For each fold \(k\), train the black-box regression \(\mhat^{(-k)}\) on the labeled observations outside \(I_k\), compute out-of-fold predictions \(\mhat^{(-k)}(X_i)\) for \(i\in I_k\), and compute predictions \(\mhat^{(-k)}(\widetilde X_j)\) on the unlabeled sample.
\item Pool the labeled out-of-fold predictions \(\{\mhat^{(-k(i))}(X_i)\}_{i=1}^n\) and fit a calibrator \(\fhat\) to the labeled pairs \((\mhat^{(-k(i))}(X_i),Y_i)\).
\item Form calibrated predictions
\[
\mtilde(X_i)=\fhat(\mhat^{(-k(i))}(X_i))
\]
and
\[
\mtilde(\widetilde X_j)=\frac{1}{K}\sum_{k=1}^K \fhat(\mhat^{(-k)}(\widetilde X_j)),
\]
and then compute the plug-in estimator \cref{eq:plugin}.
\end{enumerate}

This cross-fitted modification is the natural DML-style implementation when one seeks to reduce overfitting from estimating the initial regression on the same labeled data later used for calibration and inference \citep{schick1986asymptotically,zheng2010asymptotic,DoubleML}. Cross-calibration was proposed by \citet{van2023causal} and used for debiased inference in \citet{van2024stabilized,van2024automatic,rabenseifner2025causal}.


\section{Theory for isotonic calibration}
\label{sec:theory}
We study the isotonic-calibrated plug-in estimator \(\widehat{\psi}_{\mathrm{iso}}\). We first establish asymptotic linearity, asymptotic normality, and Wald inference.  We then study the efficiency of \(\widehat{\psi}_{\mathrm{iso}}\), showing that it is efficient in a reduced data model and giving conditions under which it is also efficient in the full model, meaning that it attains the minimum asymptotic variance among regular estimators \citep{bickel1993efficient,van2000asymptotic}. To keep the exposition light, we focus on the non-cross-fitted estimator and treat the initial black-box score \(\mhat\) as fixed. 


\subsection{Asymptotic Linearity and Inference}
We now state the main asymptotic expansion for \(\widehat{\psi}_{\mathrm{iso}}\). Throughout this subsection, let \(f_0\) denote a population isotonic calibration map for the fixed score \(\mhat\), defined by
\begin{equation}
\label{eq:pop-iso}
f_0 \in \arg\min_{f\in\mathcal F_{\mathrm{iso}}}
\E_{P_0}\bigl[\{Y-f(\mhat(X))\}^2\bigr],
\qquad
m_0 = f_0 \circ \mhat.
\end{equation}
Then \(m_0\) is the population \(L^2\)-optimal monotone transformation of \(\mhat(X)\).

\begin{assumption}[Regularity conditions for isotonic calibration]
\label{assump:isoreg}
Assume:
\begin{enumerate}[leftmargin=6mm,label=(\roman*)]
\item \textbf{Boundedness:} there exists \(C_0<\infty\) such that \(|\E[Y \mid \mhat(X)]|\le C_0\) almost surely and \(\sup_x |m_{n,\mathrm{iso}}^\star(x)| = O_P(1)\).
\item \textbf{Controlled labeled errors:} \(Y-\E[Y\mid \mhat(X)]\) is sub-Gaussian or subexponential under \(P_0\).
\end{enumerate}
\end{assumption}

The target of the asymptotic analysis is the population isotonic limit \(m_0\), which need not equal the true regression \(\muzero\). \Cref{assump:isoreg} ensures that the classical isotonic rate \(\|m_{n,\mathrm{iso}}^\star-m_0\|_{2,\Px}^2=O_p(n^{-2/3})\) is achieved; see, for example, \citet{chatterjee2015risk,yang2018contraction}.

\begin{theorem}[Asymptotic linearity under the two-sample model]
\label{thm:al}
Suppose \cref{assump:design,assump:isoreg} hold and the calibration score equations contain the constant score \(h\equiv 1\). Then
\begin{align}
\label{eq:al}
\widehat{\psi}_{\mathrm{iso}}-\psi_0
&=
\rhon(\Pl-P_0)\Bigl\{m_0-\psi_0+\rhon^{-1}(Y-m_0)\Bigr\}
\nonumber\\
&\quad +
(1-\rhon)(\Pu-\Px)(m_0-\psi_0)
+R_{n,N},
\end{align}
where  $R_{n,N} = o_p(M^{-1/2})$ satisfies 
\[
R_{n,N}
=
O_p(n^{-2/3})
+
O_p\bigl(n^{-1/3}N^{-1/2}\bigr)
\]
Thus, under the \(\sqrt{M}\) scaling, the isotonic calibrated plug-in estimator is asymptotically linear with influence pair \((D_{m_0}^L,D_{m_0}^U)\) given by
\[
D_{m_0}^L(X,Y)=m_0(X)-\psi_0+\rhozero^{-1}\{Y-m_0(X)\},
\qquad
D_{m_0}^U(\widetilde X)=m_0(\widetilde X)-\psi_0.
\]
\end{theorem}
When \(N>n\), the finite linear approximation error satisfies \(R_{n,N}=O_p(n^{-2/3})\), where this slower rate arises from the nonparametric complexity of the isotonic calibrator \(f_n\); under linear calibration, the corresponding remainder is typically \(O_p(n^{-1})\). 

\begin{corollary}[Asymptotic normality and Wald inference]
\label{cor:normal}
Under the conditions of \cref{thm:al},
\[
\sqrt{M}\,(\widehat{\psi}_{\mathrm{iso}}-\psi_0)\overset{d}{\longrightarrow}N(0,\sigma_0^2),
\]
where
\[
\sigma_0^2
:=
\rhozero\,\Var\left\{m_0(X)-\psi_0+\rhozero^{-1}(Y-m_0(X))\right\}
+
(1-\rhozero)\,\Var\left\{m_0(\widetilde X)-\psi_0\right\}.
\]
If
\[
\widehat D_i^L := m_{n,\mathrm{iso}}^\star(X_i)-\widehat{\psi}_{\mathrm{iso}}+\rhon^{-1}\{Y_i-m_{n,\mathrm{iso}}^\star(X_i)\},
\qquad
\widehat D_j^U := m_{n,\mathrm{iso}}^\star(\widetilde X_j)-\widehat{\psi}_{\mathrm{iso}},
\]
and
\[
\widehat \sigma^2
:=
\rhon \Pl\bigl[(\widehat D^L)^2\bigr]
+
(1-\rhon)\Pu\bigl[(\widehat D^U)^2\bigr],
\]
then \(\widehat\sigma^2 \to_p \sigma_0^2\), and standard Wald confidence intervals follow.
\end{corollary}

The above theorem is largely a special case of the calibrated DML analysis in \cite{van2024automatic}. The proof combines the AIPW representation in Theorem~\ref{thm:repr}, the finite uniform entropy integral of the isotonic calibrator class to control empirical-process remainders, and standard DML/TMLE arguments \citep{bickel1993efficient,vanderlaanunified,bang2005doubly,vanderLaanRose2011}; see, for example, \citet{kennedy2024semiparametric} for a review.

\subsection{Constructing confidence intervals via the bootstrap}

 

Confidence intervals may be constructed using the standard Wald approximation from Corollary \ref{cor:normal}, namely
\[
\widehat{\psi}_{\mathrm{iso}}
\pm
z_{1-\alpha/2}\,\frac{\hat\sigma}{\sqrt{n+N}}.
\]
As an alternative, one may use the nonparametric bootstrap and refit the isotonic calibration step within each bootstrap replicate, as proposed and theoretically studied by \citet{van2024automatic}; see also \citet{tang2024consistency}. This avoids direct estimation of the influence-function variance and automatically propagates uncertainty from estimation of the calibration map.  \texttt{Python} code for both approaches is provided in Appendix \ref{appendix:code}.


For each bootstrap replicate \(b=1,\dots,B\), we sample with replacement from the labeled observations \(\{(X_i,Y_i)\}_{i=1}^n\) and, independently, from the unlabeled observations \(\{\widetilde X_j\}_{j=1}^N\), obtaining bootstrap samples
\[
\{(X_i^{*(b)},Y_i^{*(b)})\}_{i=1}^n,
\qquad
\{\widetilde X_j^{*(b)}\}_{j=1}^N.
\]
We then refit the isotonic calibrator on the bootstrap labeled sample:
\begin{align*}
\hat f^{*(b)}
\in
\arg\min_{f \in \mathcal F_{\mathrm{iso}}}
\sum_{i=1}^n
\bigl\{Y_i^{*(b)} - f(\mhat(X_i^{*(b)}))\bigr\}^2,
\end{align*}
and define
\[
m_{\mathrm{iso}}^{*(b)}(X):=\hat f^{*(b)}\{\mhat(X)\}.
\]
The corresponding bootstrap replicate of the estimator is
\[
\widehat{\psi}_{\mathrm{iso}}^{*(b)}
:=
\frac{1}{n+N}\left\{
\sum_{i=1}^n m_{\mathrm{iso}}^{*(b)}(X_i^{*(b)})
+
\sum_{j=1}^N m_{\mathrm{iso}}^{*(b)}(\widetilde X_j^{*(b)})
\right\}.
\]

After computing \(B\) bootstrap replicates, one may form a percentile interval using the empirical \(\alpha/2\) and \(1-\alpha/2\) quantiles of
\[
\bigl\{\widehat{\psi}_{\mathrm{iso}}^{*(b)}\bigr\}_{b=1}^B.
\]
Alternatively, one may use the bootstrap standard error
\[
\widehat{\sigma}_{\mathrm{boot}}^2
:=
\frac{1}{B-1}\sum_{b=1}^B
\left(\widehat{\psi}_{\mathrm{iso}}^{*(b)}-\bar{\psi}_{\mathrm{iso}}^*\right)^2,
\qquad
\bar{\psi}_{\mathrm{iso}}^*
:=
\frac{1}{B}\sum_{b=1}^B \widehat{\psi}_{\mathrm{iso}}^{*(b)},
\]
and report the normal-approximation interval
\[
\widehat{\psi}_{\mathrm{iso}}
\pm
z_{1-\alpha/2}\,\widehat{\sigma}_{\mathrm{boot}}.
\]

The key implementation detail is that the isotonic regression step must be refit within each bootstrap replicate rather than held fixed. Otherwise, the resulting interval does not reflect uncertainty from estimation of the calibration map.


\subsection{Efficiency considerations}

In this section, we clarify the efficiency properties of the isotonic-calibrated plug-in estimator \(\widehat{\psi}_{\mathrm{iso}}\). The asymptotic linearity result of Theorem~\ref{thm:al} yields the influence pair
\[
D_{m_0}^L(X,Y)=m_0(X)-\psi_0+\rhozero^{-1}\{Y-m_0(X)\},
\qquad
D_{m_0}^U(\widetilde X)=m_0(\widetilde X)-\psi_0.
\]
The key point is that \(\widehat{\psi}_{\mathrm{iso}}\) uses the covariates only through the one-dimensional score \(m_0(X)\). It is therefore natural to assess its efficiency relative to a model in which \(m_0(X)\), rather than the full covariate \(X\), is observed. In that reduced two-sample model, calibration implies \(\E[m_0(X)]=\E[Y]=\psi_0\), and the pair above is the efficient influence pair for \(\psi_0\). We then compare this reduced-model benchmark with the full-model benchmark and identify conditions under which they coincide. Our semiparametric efficiency results for the two-sample setting are straightforward applications of more general data-fusion results in \citet{li2023efficient,graham2024towards}.

\begin{proposition}[Efficient influence pair in the reduced model indexed by \(m_0\)]
\label{prop:eif-red}
Consider the reduced two-sample experiment
\[
\mathcal E_{m_0}
:=
\Bigl(
\{(m_0(X_i),Y_i)\}_{i=1}^n,\,
\{m_0(\widetilde X_j)\}_{j=1}^N
\Bigr),
\]
obtained by replacing the full covariate \(X\) with the scalar summary \(m_0(X)\). On the submodel satisfying
\[
\E_P[Y]=\E_P[m_0(X)],
\]
the efficient influence pair for \(\psi(P)=\E_P[Y]\) is
\begin{align}
\label{eq:eif-red-L}
D_{m_0}^L(X,Y)
&=
m_0(X)-\psi_0+\rhozero^{-1}\{Y-m_0(X)\},\\
\label{eq:eif-red-U}
D_{m_0}^U(\widetilde X)
&=
m_0(\widetilde X)-\psi_0.
\end{align}
\end{proposition}

\begin{corollary}[Efficiency regimes]
\label{cor:efficiency}
Under the conditions of \cref{thm:al}, the following hold.
\begin{enumerate}[leftmargin=6mm,label=(\roman*)]
\item The influence pair in \cref{thm:al} coincides with the efficient influence pair in \cref{prop:eif-red}. Hence \(\widehat{\psi}_{\mathrm{iso}}\) is efficient for \(\psi_0\) in the coarsened experiment \(\mathcal E_{m_0}\).
\item If \(m_0=\muzero\) almost surely, then the influence pair in \cref{thm:al} coincides with the efficient influence pair in \cref{prop:eif-full}. Hence \(\widehat{\psi}_{\mathrm{iso}}\) is efficient for \(\psi_0\) in the full two-sample experiment.
\item If \(\E[Y\mid m(X)]\) is monotone in \(m(X)\), then \(m_0(X)=\E[Y\mid m(X)]\) almost surely. Hence \(\widehat{\psi}_{\mathrm{iso}}\) is efficient for \(\psi_0\) in the coarsened experiment \(\mathcal E_m\).
\end{enumerate}
\end{corollary}



\begin{remark}[Practical takeaway]
From a practical perspective, isotonic calibration can improve efficiency whenever the original prediction score is informative but miscalibrated. The estimator remains a simple plug-in estimator, but it behaves as if one had replaced the original score by its best monotone recalibration.
\end{remark}

\begin{remark}[When full efficiency is attained]
Full efficiency requires that the calibrated target recover all relevant regression information. This holds if \(m_0=\mu_0\) almost surely, and more generally fails when the one-dimensional score \(m(X)\) omits predictive information in \(X\) that cannot be recovered through monotone recalibration.
\end{remark}



 

 
 

\section{Experiments}
\label{sec:exp}

\subsection{Simulation study}
 
\label{sec:simulation}
We first study a controlled binary-outcome simulation designed to isolate the effect of score miscalibration. Let \(S\sim N(0,1)\), define
\[
\mu_0(S)=\{1+\exp(-5S)\}^{-1},
\]
and draw \(Y\mid S\sim \mathrm{Bernoulli}\{\mu_0(S)\}\). We compare two score regimes. In the mildly miscalibrated regime, we first add small probability-scale noise and then apply a gentle Platt-style logit tilt,
\[
m(X)=\sigma\Bigl(1.5\,\mathrm{logit}\bigl(\mathrm{clip}\{\mu_0(S)+\varepsilon,\,0.01,\,0.99\}\bigr)\Bigr).
\]
In the poorly calibrated regime, we apply a stronger affine-biased monotone logit distortion,
\[
m(X)=\mathrm{clip}\left[-0.15+0.75\,\sigma\bigl\{0.8\,z(S)+0.1\,z(S)^3-1\bigr\},\,0.01,\,0.99\right],
\]
where \(z(S):=\mathrm{logit}\{\mu_0(S)\}\). This preserves ranking information but distorts the score on the probability scale in a genuinely nonlinear way. We report the labeled-only mean (\textsc{Naive}), PPI, AIPW, linear recalibration, and isotonic recalibration. The labeled sample size ranges over \(n\in\{50,100,200,400,800,1200,2400\}\). The main-text figure focuses on the poorly calibrated regime and shows both the balanced regime \(N/n=1\) and the large-unlabeled regime \(N/n=16\), which together illustrate how the relative behavior changes with unlabeled-sample size. The mildly miscalibrated setting is discussed briefly below and deferred to Appendix~\ref{app:simulation-extra}, where the score-based estimators are nearly indistinguishable. All summaries are based on 5000 Monte Carlo repetitions.

\begin{figure}[htb]
\centering
\includegraphics[width=\textwidth]{exp_plots_tables/sim_fig_main.pdf}
\caption{Poorly calibrated synthetic simulation study in two unlabeled-sample regimes. The top row shows the balanced case \(N/n=1\), and the bottom row shows the large-unlabeled case \(N/n=16\). The four columns show absolute bias, empirical Monte Carlo standard deviation, coverage, and relative efficiency versus PPI, computed from empirical Monte Carlo variance. AIPW and PPI are most clearly separated in the balanced regime, while the two become much closer once unlabeled data are abundant.}
\label{fig:simulation-rmse}
\end{figure}

\cref{fig:simulation-rmse} reports the four diagnostics most relevant to the theory in two representative unlabeled-sample regimes, \(N/n=1\) and \(N/n=16\): absolute bias, empirical Monte Carlo standard deviation, coverage, and relative efficiency versus PPI. In the mildly miscalibrated regime, all score-based estimators perform similarly, so recalibration offers little beyond the classical AIPW correction; we therefore relegate those curves to Appendix~\ref{app:simulation-extra}. Under stronger miscalibration, by contrast, recalibration improves efficiency by reducing the residual variance of the score-based estimator, most clearly in the empirical-variability and relative-efficiency panels. At the balanced ratio \(N/n=1\), AIPW is slightly more efficient than PPI in small samples and is comparable in larger samples, whereas at \(N/n=16\) it can become slightly less efficient under substantial misspecification in \(m\). LinearCal captures much of this gain already in the balanced regime, while IsoCal can improve further once enough labeled data are available to support a stable nonlinear correction, especially when unlabeled data are abundant. Coverage remains broadly reasonable throughout, although the more aggressive nonlinear recalibrators can mildly undercover in the larger-sample setting.  


\subsection{Empirical Illustration}
\label{sec:empirical}

We evaluate the proposed calibration-based estimators on the mean-estimation benchmarks used in the original experiments of \citet{ppi2023,angelopoulos2023ppi}, as provided in the \texttt{ppi\_py} package. Our goal is not to revisit the original benchmark comparison, but rather to understand when one-dimensional recalibration of a prediction score improves semisupervised mean inference. We use the same \texttt{ppi\_py} datasets and retain the labeled-sample-size grids from the original examples. For each labeled sample size \(n\), we repeatedly split the fully labeled benchmark into a labeled subset of size \(n\) and an unlabeled subset of size \(N\), apply all estimators to the same split, and treat the full-sample mean as the ground truth.\footnote{Because repeated sample splits of a fixed finite benchmark dataset are not i.i.d.\ draws from a superpopulation, asymptotic Wald confidence intervals can exhibit either undercoverage or overcoverage. This is not specific to our calibration extensions, but is a feature of evaluating Wald-type intervals under finite-population resampling.} Results are averaged over 5000 random splits for each \(n\).



\paragraph{Benchmarks and baselines.}
We study the \texttt{forest}, \texttt{galaxies}, and \texttt{census\_income} semisupervised mean-estimation examples from the official benchmark suite. The retained baselines are the original labeled-only interval, the imputation benchmark on the binary-outcome datasets, the semisupervised linear-regression baseline from \texttt{ppi\_py} that uses the full covariate vector, PPI, and the classical AIPW estimator based on the same prediction score. The imputation benchmark averages model predictions over the pooled sample without a labeled residual correction, whereas the semisupervised linear baseline fits a linear regression on the full covariates and combines unlabeled prediction averaging with a labeled residual correction. AIPW is the corresponding pooled plug-in plus augmentation estimator based on the original score, whereas PPI is a less efficient variant that discards labeled covariate information in the plug-in term.

We additionally consider three recalibration-based comparators that fall directly within our framework: linear recalibration (\texttt{LinearCal}), Platt scaling (\texttt{Platt}) on the binary-outcome datasets, and isotonic recalibration with a fixed minimum of 10 observations per bin (\texttt{IsoCal}). For the binary-outcome datasets, we also report the Venn--Abers shrinkage variant described in Section~\ref{sec:iso}. Platt scaling and Venn--Abers are omitted on \texttt{census\_income} because the outcome is not binary. We do not include PPI++ as separate baselines, since they are algebraically equivalent to linear recalibration/prognostic-score adjustment. To keep the displayed figures and appendix tables readable, we omit the semisupervised linear baseline from those summaries and focus there on labeled-only, imputation, PPI, AIPW, and the calibration-based estimators.

\paragraph{Metrics and main findings.}
For each estimator, we report Monte Carlo bias, empirical variance, MSE, Wald-interval coverage, and relative efficiency versus PPI, where relative efficiency is computed from empirical Monte Carlo variances. Estimated standard errors are used only to construct Wald intervals; all reported variability summaries are empirical Monte Carlo quantities. 

The results suggest a simple practical picture. AIPW is the natural estimator based on the original score, and it typically matches or outperforms PPI. In \texttt{forest}, where the original score is already strong, PPI becomes less efficient as the labeled sample grows, whereas in \texttt{galaxies} the two methods remain nearly tied. When the unlabeled sample is much larger than the labeled sample, however, AIPW and PPI become nearly indistinguishable because the labeled contribution to the pooled plug-in term is negligible; this is what we observe in \texttt{census\_income}, where the unlabeled sample is enormous.

Recalibration is most useful when the score is meaningfully miscalibrated and the labeled sample is large enough to estimate the second-stage calibrator stably. In that regime, linear recalibration often captures most of the available gain under mild misspecification, as in \texttt{census\_income}, whereas more flexible monotone calibration is mainly competitive once the labeled sample is moderately large, with \texttt{galaxies} providing the clearest example. Overall, these results leave little practical motivation for PPI: when \(m\) is accurate, AIPW is more efficient, and when \(m\) is inaccurate, even simple recalibration methods such as LinearCal often remove much of the misspecification bias.
 
\begin{figure}[htb]
\centering
\includegraphics[width=\textwidth]{exp_plots_tables/fig_paper_main_grid.pdf}
\caption{Main-text benchmark summary showing MSE, relative efficiency versus PPI, and coverage across the reproduced PPI benchmarks, restricted to the most informative score-based comparators: AIPW, PPI, LinearCal, and IsoCal (minimum bin size 10). Relative efficiency is computed from empirical Monte Carlo variance, so values above one indicate improvement over PPI. The dashed horizontal line in the coverage panels marks the nominal target coverage, which is \(0.95\) for \texttt{forest} and \(0.90\) for \texttt{galaxies} and \texttt{census\_income}.}
\label{fig:ppi-main}
\end{figure}

\begin{table}[htb]
\centering
\small
\begin{tabular}{p{0.16\textwidth}p{0.34\textwidth}p{0.34\textwidth}}
\toprule
Dataset & Smaller labeled-sample regime & Larger labeled-sample regime \\
\midrule
\texttt{forest} & AIPW weakly dominates PPI; the added flexibility of LinearCal and IsoCal leads to higher finite-sample variance relative to AIPW. & AIPW is substantially more efficient than PPI, while additional recalibration yields little beyond modest MSE gains. \\
\texttt{galaxies} & AIPW and PPI are very close; LinearCal is the closest recalibration competitor, and IsoCal also remains competitive. & The leading score-based methods become very similar, so simple corrections are usually sufficient. \\
\texttt{census\_income} & PPI and AIPW are effectively indistinguishable because the unlabeled sample is very large; flexible calibration can be too aggressive when the labeled sample is very small. & LinearCal performs best on average, with IsoCal also competitive once the labeled sample is large enough. \\
\bottomrule
\end{tabular}
\caption{Compact summary of the reproduced PPI benchmarks across regimes. The main empirical distinction is between small labeled-sample regimes, where simple or no recalibration is safer, and larger labeled-sample regimes, where recalibration can help when the score is miscalibrated.}
\label{tab:ppi-takeaways}
\end{table}

 


\section{Conclusion}

We studied semisupervised mean estimation with a black-box prediction score and proposed a simple calibrated plug-in procedure: fit the score, recalibrate it on the labeled sample, and average the recalibrated predictions over the pooled covariate sample. Our main observation is that this seemingly simple plug-in rule also admits an exact AIPW representation, placing it directly within the standard semiparametric AIPW framework. In particular, isotonic calibration yields a lightweight, tuning-free estimator with asymptotic linearity, asymptotic normality, and valid Wald inference under weak conditions. Wald-type confidence intervals may be constructed from the influence function, or equivalently from the residual variance. Alternatively, the bootstrap may be used to better capture the finite-sample variability of the isotonic regression update.

From an efficiency perspective, the isotonic-calibrated estimator is efficient for the information retained by the calibrated score \(m_0(X)\). It is nonparametrically efficient for the full covariate information in \(X\) when the calibrated limit \(m_0\) coincides with the true regression function \(\mu_0\). Isotonic calibration always improves efficiency over the simple AIPW estimator based on the original score \(m(X)\), and it typically also improves on linear calibration. Empirically, the results suggest a simple practical message: when the original score is already well calibrated, recalibration changes little, whereas under meaningful miscalibration, even simple affine or monotone recalibration can improve semisupervised efficiency while preserving a transparent plug-in form. We also showed that PPI++ is equivalent to AIPW with empirical efficiency maximization \citep{rubin2008empirical}, and can equivalently be understood as linear calibration, classical prognostic-score regression adjustment \citep{hansen2008prognostic,rosenblum2009using,hojbjerre2025powering}, and a special case of TMLE \citep{van2006targeted}.

While we focus on calibration of regression functions and inference on means, the ideas here may extend to more general loss-based supervised learning tasks, such as median and quantile regression \citep{noarov2023statistical,jung2021moment,roth2022uncertain,whitehouse2024orthogonal}. In particular, when the target parameter is defined by a moment equation, as in \(M\)-estimation, one could debias or calibrate the corresponding augmented moment equation; see Appendix \ref{appendix::momentgeneral}. However, this would generally require recalibration at each candidate parameter value, which is computationally impractical and may lead to difficult-to-optimize estimating equations, including nonconvex \(M\)-estimation problems. A more practical alternative is a sieve-based, TMLE-style multiaccuracy adjustment of the predictor or weights, as in the Efficient Plug-in Learning framework of \citet{van2024combining}, which addresses this issue for nonconvex Neyman-orthogonal loss functions \citep{foster2023orthogonal}. Inference for generic smooth functionals of an \(M\)-estimand is studied by \citet{van2025automatic}.

\bibliographystyle{plainnat}
\bibliography{references}

 

 


\appendix
 
\section{Equivalence between PPI and PPI++ and linear calibration}
\label{sec::ppi++}

\begin{proposition}[PPI++ is algebraically equivalent to affine linear calibration]
\label{prop:ppi++-linear}
Let \(m_i:=m(X_i)\) for \(i=1,\dots,n\) and \(\tilde m_j:=m(\widetilde X_j)\) for \(j=1,\dots,N\). Define the linearly calibrated fit by
\[
(\hat a_{\mathrm{lin}},\hat b_{\mathrm{lin}})
\in
\arg\min_{a,b\in\R}\sum_{i=1}^n \{Y_i-a\,m_i-b\}^2,
\]
and the associated plug-in estimator by
\[
\widehat\psi_{\mathrm{lin}}
=
\rho_n \Pl\{\hat a_{\mathrm{lin}} m(X)+\hat b_{\mathrm{lin}}\}
+
(1-\rho_n)\Pu\{\hat a_{\mathrm{lin}} m(\widetilde X)+\hat b_{\mathrm{lin}}\}.
\]

Now define the PPI++ estimator as the AIPW estimator over the scaling class
\(\{\lambda m(X):\lambda\in\R\}\),
\[
\widehat\psi_{++}(\lambda)
:=
\rho_n \Pl\{\lambda m(X)\}
+
(1-\rho_n)\Pu\{\lambda m(\widetilde X)\}
+
\Pl\{Y-\lambda m(X)\},
\]
evaluated at the empirical efficiency-maximizing coefficient \citep{rubin2008empirical,angelopoulos2023ppi}
\[
\hat\lambda_{++}
=
\frac{\Pl\!\bigl[(m(X)-\Pl m(X))(Y-\Pl Y)\bigr]}
{\Pl\!\bigl[(m(X)-\Pl m(X))^2\bigr]},
\]
with the convention \(\hat\lambda_{++}=0\) if the denominator is zero.

Then
\[
\hat\lambda_{++}=\hat a_{\mathrm{lin}}
\qquad\text{and}\qquad
\widehat\psi_{++}(\hat\lambda_{++})=\widehat\psi_{\mathrm{lin}}.
\]
Hence PPI++ and affine linear calibration are algebraically identical in the one-score mean-estimation problem.
\end{proposition}

\begin{proof}
The normal equations for the linear regression fit are
\[
\Pl\{Y-\hat a_{\mathrm{lin}}m(X)-\hat b_{\mathrm{lin}}\}=0
\]
and
\[
\Pl\bigl[m(X)\{Y-\hat a_{\mathrm{lin}}m(X)-\hat b_{\mathrm{lin}}\}\bigr]=0.
\]
The first equation gives
\[
\hat b_{\mathrm{lin}}=\Pl Y-\hat a_{\mathrm{lin}}\Pl m(X).
\]
Substituting this into the second equation yields
\[
\Pl\bigl[(m(X)-\Pl m(X))(Y-\Pl Y)\bigr]
=
\hat a_{\mathrm{lin}}\,
\Pl\bigl[(m(X)-\Pl m(X))^2\bigr].
\]
Therefore
\[
\hat a_{\mathrm{lin}}
=
\frac{\Pl\!\bigl[(m(X)-\Pl m(X))(Y-\Pl Y)\bigr]}
{\Pl\!\bigl[(m(X)-\Pl m(X))^2\bigr]}
=
\hat\lambda_{++},
\]
with the same zero-variance convention if \(\Pl[(m-\Pl m)^2]=0\).

It remains to compare the two estimators. Using
\(\hat b_{\mathrm{lin}}=\Pl Y-\hat a_{\mathrm{lin}}\Pl m(X)\),
we have
\begin{align*}
\widehat\psi_{\mathrm{lin}}
&=
\rho_n \Pl\{\hat a_{\mathrm{lin}}m(X)+\hat b_{\mathrm{lin}}\}
+
(1-\rho_n)\Pu\{\hat a_{\mathrm{lin}}m(\widetilde X)+\hat b_{\mathrm{lin}}\}\\
&=
\hat b_{\mathrm{lin}}
+
\hat a_{\mathrm{lin}}
\Bigl[
\rho_n \Pl m(X)+(1-\rho_n)\Pu m(\widetilde X)
\Bigr]\\
&=
\Pl Y-\hat a_{\mathrm{lin}}\Pl m(X)
+
\hat a_{\mathrm{lin}}
\Bigl[
\rho_n \Pl m(X)+(1-\rho_n)\Pu m(\widetilde X)
\Bigr]\\
&=
\Pl Y
+
(1-\rho_n)\hat a_{\mathrm{lin}}
\bigl\{\Pu m(\widetilde X)-\Pl m(X)\bigr\}.
\end{align*}
On the other hand,
\begin{align*}
\widehat\psi_{++}(\hat\lambda_{++})
&=
\rho_n \Pl\{\hat\lambda_{++} m(X)\}
+
(1-\rho_n)\Pu\{\hat\lambda_{++} m(\widetilde X)\}
+
\Pl\{Y-\hat\lambda_{++} m(X)\}\\
&=
\Pl Y
+
(1-\rho_n)\hat\lambda_{++}
\bigl\{\Pu m(\widetilde X)-\Pl m(X)\bigr\}.
\end{align*}
Since \(\hat\lambda_{++}=\hat a_{\mathrm{lin}}\), the two estimators are equal:
\[
\widehat\psi_{++}(\hat\lambda_{++})=\widehat\psi_{\mathrm{lin}}.
\]
This proves the claim.
\end{proof}


\begin{proposition}[PPI and AIPW as intercept-only calibration]
\label{prop:ppi-aipw-const}
Let
\[
\hat a_{\mathrm{const}}
\in
\arg\min_{a\in\R}\sum_{i=1}^n \{Y_i-a-\mhat(X_i)\}^2,
\qquad
m_{n,\mathrm{const}}^\star(X):=\hat a_{\mathrm{const}}+\mhat(X).
\]
Then
\[
\hat a_{\mathrm{const}}=\Pl\{Y-\mhat(X)\}.
\]
Consequently, the PPI estimator
\[
\widehat{\psi}_{\mathrm{PPI}}
:=
\Pu\{\mhat(\widetilde X)\}+\Pl\{Y-\mhat(X)\}
\]
satisfies
\[
\widehat{\psi}_{\mathrm{PPI}}
=
\Pu\{m_{n,\mathrm{const}}^\star(\widetilde X)\},
\]
while the standard AIPW estimator
\[
\widehat{\psi}_{\mathrm{AIPW}}
:=
\rho_n \Pl\{\mhat(X)\}
+
(1-\rho_n)\Pu\{\mhat(\widetilde X)\}
+
\Pl\{Y-\mhat(X)\}
\]
satisfies
\[
\widehat{\psi}_{\mathrm{AIPW}}
=
\rho_n \Pl\{m_{n,\mathrm{const}}^\star(X)\}
+
(1-\rho_n)\Pu\{m_{n,\mathrm{const}}^\star(\widetilde X)\}.
\]
Hence PPI is the unlabeled-only plug-in estimator based on the mean-calibrated score \(m_{n,\mathrm{const}}^\star\), whereas AIPW is the corresponding pooled plug-in estimator.
\end{proposition}

\begin{proof}
The objective in the definition of \(\hat a_{\mathrm{const}}\) is quadratic in \(a\), and its first-order condition is
\[
0
=
\Pl\{Y-a-\mhat(X)\}.
\]
Therefore
\[
\hat a_{\mathrm{const}}=\Pl\{Y-\mhat(X)\}.
\]

Using \(m_{n,\mathrm{const}}^\star(\widetilde X)=\mhat(\widetilde X)+\hat a_{\mathrm{const}}\), we obtain
\[
\Pu\{m_{n,\mathrm{const}}^\star(\widetilde X)\}
=
\Pu\{\mhat(\widetilde X)\}+\hat a_{\mathrm{const}}
=
\Pu\{\mhat(\widetilde X)\}+\Pl\{Y-\mhat(X)\}
=
\widehat{\psi}_{\mathrm{PPI}}.
\]

Likewise,
\begin{align*}
\rho_n \Pl\{m_{n,\mathrm{const}}^\star(X)\}
+(1-\rho_n)\Pu\{m_{n,\mathrm{const}}^\star(\widetilde X)\}
&=
\rho_n \Pl\{\mhat(X)+\hat a_{\mathrm{const}}\}
+(1-\rho_n)\Pu\{\mhat(\widetilde X)+\hat a_{\mathrm{const}}\}\\
&=
\rho_n \Pl\{\mhat(X)\}
+(1-\rho_n)\Pu\{\mhat(\widetilde X)\}
+\hat a_{\mathrm{const}}\\
&=
\rho_n \Pl\{\mhat(X)\}
+(1-\rho_n)\Pu\{\mhat(\widetilde X)\}
+\Pl\{Y-\mhat(X)\}\\
&=
\widehat{\psi}_{\mathrm{AIPW}}.
\end{align*}
This proves the claim.
\end{proof}




\section{Prediction-Powered Causal Inference via Calibration (i.e., Calibrated DML)}
\label{app:causal}
This appendix records the direct translation of the two-sample missing-outcome setup to randomized experiments. The central point is that the semisupervised problem studied in the main text and two-arm causal inference in a randomized trial are the same statistical problem once the unobserved potential outcome in one arm is interpreted as the missing outcome. We then illustrate how our calibrated plug-in estimators apply in this setting, where they arise as special cases of calibrated DML \citep{van2024automatic}.

We use the standard potential-outcomes notation. Each unit has two potential outcomes, \(Y(1)\) under treatment and \(Y(0)\) under control, and the observed outcome satisfies
\[
Y=AY(1)+(1-A)Y(0),
\]
where \(A\in\{0,1\}\) is the treatment assignment. The mean potential outcomes are
\[
\mu_1:=\E\{Y(1)\},
\qquad
\mu_0:=\E\{Y(0)\},
\]
and the average treatment effect is \(\tau:=\mu_1-\mu_0\). In a randomized trial, identification follows from the usual causal conditions: consistency and no interference, exchangeability induced by randomization,
\[
A\perp (X,Y(1),Y(0)),
\]
and positivity, so that each arm has positive assignment probability.

Now partition the observed data into the two arm-specific samples
\[
\mathcal D_{n,1}:=\{(X_i,Y_i):A_i=1\},
\qquad
\mathcal D_{n,0}:=\{(\widetilde X_j,\widetilde Y_j):A_j=0\},
\]
with sample sizes \(n_1:=|\mathcal D_{n,1}|\), \(n_0:=|\mathcal D_{n,0}|\), and \(n=n_1+n_0\). Let \(\mathbb P_{n,a}\) denote the empirical mean over arm \(a\), and let \(\hat\pi_a:=n_a/n\). To estimate \(\mu_1\), one treats \(\mathcal D_{n,1}\) as the labeled sample for \(Y(1)\) and uses the covariates in \(\mathcal D_{n,0}\) as the unlabeled sample. To estimate \(\mu_0\), one reverses the roles of the two arms. Thus each mean potential outcome is a two-sample missing-outcome estimand of exactly the form analyzed in the main text.

For \(a\in\{0,1\}\), let \(m_a(X)\) be an arm-specific prediction score fit using only observations in arm \(a\). The standard arm-specific AIPW estimator is
\[
\hat\mu_a^{\mathrm{AIPW}}
:=
\mathbb P_n\{m_a(X)\}
+
\mathbb P_n\!\left[
\frac{\1(A=a)}{\hat\pi_a}\{Y-m_a(X)\}
\right].
\]


For \(a\in\{0,1\}\), let \(m_{n,a}^*(X)\) denote the recalibrated version of the score \(m_a(X)\) The calibrated estimator of \(\mu_a\) is then
\begin{equation}
\label{eq:rct-mua-plugin}
\hat\mu_a^{\mathrm{cal}}
:=
\mathbb P_n\{m_{n,a}^*(X)\}
=
\hat\pi_a\,\mathbb P_{n,a}\{m_{n,a}^*(X)\}
+
(1-\hat\pi_a)\,\mathbb P_{n,1-a}\{m_{n,a}^*(X)\}.
\end{equation}
This is exactly the pooled plug-in estimator from the two-sample formulation, with one arm contributing outcomes and the other contributing only covariates. The same estimator also has the exact AIPW form
\begin{equation}
\label{eq:rct-mua-aipw}
\hat\mu_a^{\mathrm{cal}}
=
\mathbb P_n\{m_{n,a}^*(X)\}
+
\mathbb P_n\!\left[
\frac{\1(A=a)}{\hat\pi_a}\{Y-m_{n,a}^*(X)\}
\right].
\end{equation}
Because
\[
\mathbb P_n\!\left[
\frac{\1(A=a)}{\hat\pi_a}\{Y-m_{n,a}^*(X)\}
\right]
=
\mathbb P_{n,a}\{Y-m_{n,a}^*(X)\},
\]
the residual correction vanishes exactly whenever the calibrator includes the constant score and therefore enforces mean calibration within arm \(a\). This is the randomized-trial analogue of the exact augmented representation in \cref{thm:repr}.


The average treatment effect is then estimated by
\begin{equation}
\label{eq:rct-ate}
\hat\tau^{\mathrm{cal}}
:=
\hat\mu_1^{\mathrm{cal}}-\hat\mu_0^{\mathrm{cal}}.
\end{equation}
Equivalently,
\[
\hat\tau^{\mathrm{cal}}
=
\mathbb P_n\{\tilde m_1(X)-\tilde m_0(X)\}
+
\mathbb P_n\!\left[
\frac{A}{\hat\pi_1}\{Y-\tilde m_1(X)\}
-
\frac{1-A}{\hat\pi_0}\{Y-\tilde m_0(X)\}
\right].
\]
Its influence function is the difference of the two arm-specific AIPW influences,
\[
D_\tau(O)
=
\left\{\tilde m_1(X)-\mu_1+\frac{A}{\pi_1}\bigl(Y-\tilde m_1(X)\bigr)\right\}
-
\left\{\tilde m_0(X)-\mu_0+\frac{1-A}{\pi_0}\bigl(Y-\tilde m_0(X)\bigr)\right\}.
\]
Accordingly, Wald inference for the ATE may be based on the empirical variance of the estimated influence values divided by \(n\), or on the nonparametric bootstrap refitting the calibration step within each bootstrap sample.

The practical implementation is therefore immediate. One splits the trial into \(\mathcal D_{n,1}\) and \(\mathcal D_{n,0}\), treats one arm as labeled and the other as unlabeled when estimating a given mean potential outcome, fits and recalibrates an arm-specific score, averages the recalibrated predictions over the pooled covariate sample, and finally subtracts the two arm-specific estimates to obtain the ATE. In this way, the semisupervised mean estimators in the main text become direct estimators of mean potential outcomes and treatment effects in randomized experiments.  


\section{General Moment Equations Under MAR}
\label{appendix::momentgeneral}
We briefly review the general AIPW class for moment-equation targets under missing at random; see, for example, \citet{robins1995analysis,vanderlaanunified}. Let
\[
O=(X,R,RY),
\]
where \(R\in\{0,1\}\) indicates whether \(Y\) is observed, and suppose
\[
R \perp Y \mid X,
\qquad
\pi_0(X):=\Prob(R=1\mid X)>0 \ \text{a.s.}
\]
Let the target \(\theta_0\in\Theta\subseteq\R^d\) be identified by
\[
\E[\varphi(Y,\theta_0)]=0,
\]
where \(\varphi:\mathcal Y\times\Theta\to\R^d\) is a known identifying function. Define
\[
A_0:=\partial_\theta \E[\varphi(Y,\theta)]\big|_{\theta=\theta_0},
\]
and assume \(A_0\) is nonsingular.

\begin{theorem}[AIPW class for general moment equations under MAR]
\label{thm:general-moment-aipw}
For any square-integrable function \(m:\mathcal X\times\Theta\to\R^d\), define
\begin{equation}
\label{eq:general-aipw-eq}
\psi(O;\theta,m,\pi_0)
:=
m(X,\theta)+\frac{R}{\pi_0(X)}\{\varphi(Y,\theta)-m(X,\theta)\}.
\end{equation}
Then
\[
\E[\psi(O;\theta_0,m,\pi_0)]=0.
\]
Hence any estimator \(\hat\theta\) solving the empirical estimating equation
\[
\mathbb P_n \psi(O;\hat\theta,m,\pi_0)=0
\]
is a valid AIPW-style estimator of \(\theta_0\).

Moreover, the corresponding influence function is
\begin{equation}
\label{eq:general-if}
D_m(O)
=
-A_0^{-1}\psi(O;\theta_0,m,\pi_0)
=
-A_0^{-1}\left[
m(X,\theta_0)+\frac{R}{\pi_0(X)}\{\varphi(Y,\theta_0)-m(X,\theta_0)\}
\right].
\end{equation}
Thus the class of regular influence-function-based estimators is indexed by the choice of augmentation function \(m\).

The efficient choice in the full MAR model is
\begin{equation}
\label{eq:general-efficient-m}
m_0(X,\theta_0):=\E[\varphi(Y,\theta_0)\mid X],
\end{equation}
which yields the efficient influence function
\begin{equation}
\label{eq:general-eif}
D_{\mathrm{eff}}(O)
=
-A_0^{-1}\left[
m_0(X,\theta_0)+\frac{R}{\pi_0(X)}\{\varphi(Y,\theta_0)-m_0(X,\theta_0)\}
\right].
\end{equation}
\end{theorem}


\begin{corollary}[Two-sample/PPI form]
Suppose the labeled and unlabeled samples arise from the semisupervised model with constant labeling probability \(\rho_0\), and let \(S=f(X)\) be a chosen score. Replacing \(X\) by \(S\) gives the reduced AIPW class
\[
\psi(O;\theta,m,\rho_0)
=
m(S,\theta)+\frac{R}{\rho_0}\{\varphi(Y,\theta)-m(S,\theta)\}.
\]
Equivalently, the sample estimating equation can be written as
\[
0=
\rho_n \mathbb P_n^L\{m(S,\theta)\}
+
(1-\rho_n)\mathbb P_N^U\{m(S,\theta)\}
+
\mathbb P_n^L\{\varphi(Y,\theta)-m(S,\theta)\}.
\]
The efficient choice within the reduced score model is
\[
m_0(S,\theta_0)=\E[\varphi(Y,\theta_0)\mid S].
\]
\end{corollary}





\section{Proof of the representation theorem}
\label{app:repr}

\begin{proof}[Proof of \cref{thm:repr}]
Because \cref{eq:cal-score} holds for every \(h\in\mathcal F\), any empirical balancing weight \(\what(\mtilde(X))\in\mathcal F\) satisfies
\[
\Pl\bigl[\what(\mtilde(X))\{Y-\mtilde(X)\}\bigr]=0.
\]
Adding this zero term to the plug-in estimator \cref{eq:plugin} yields \cref{eq:abw-general}.

For the second claim, since the identity map belongs to \(\mathcal F\), taking \(f(t)=t\) in \cref{eq:F-balance} gives
\[
\Pl\bigl[\what_n^*(\mtilde(X))\,\mtilde(X)\bigr]
=
\rhon \Pl\{\mtilde(X)\}
+
(1-\rhon)\Pu\{\mtilde(\widetilde X)\}
=
\psihat.
\]
Since \(\what_n^*\in\mathcal F\), \cref{eq:cal-score} also gives
\[
\Pl\bigl[\what_n^*(\mtilde(X))\{Y-\mtilde(X)\}\bigr]=0.
\]
Therefore,
\[
\Pl\bigl[\what_n^*(\mtilde(X))Y\bigr]
=
\Pl\bigl[\what_n^*(\mtilde(X))\,\mtilde(X)\bigr]
+
\Pl\bigl[\what_n^*(\mtilde(X))\{Y-\mtilde(X)\}\bigr]
=
\psihat,
\]
which is \cref{eq:abw-general-balance}.
\end{proof}


\begin{lemma}[Blockwise residual orthogonality]
\label{lem:block}
Let \(B_1,\dots,B_J\) be the pooled adjacent violator blocks of the isotonic fit, and let \(m_{n,\mathrm{iso}}^\star(X_i)=c_j\) for \(i \in B_j\). Then
\[
\sum_{i \in B_j}(Y_i-c_j)=0,
\qquad j=1,\dots,J.
\]
Consequently, for every measurable function \(h:\R \to \R\) such that \(h(m_{n,\mathrm{iso}}^\star(X_i))\) is finite for all \(i\),
\[
\sum_{i=1}^n h(m_{n,\mathrm{iso}}^\star(X_i))\{Y_i-m_{n,\mathrm{iso}}^\star(X_i)\}=0.
\]
\end{lemma}

\begin{proof}
The fitted values from isotonic regression are constant on pooled adjacent violator blocks. The KKT conditions imply that residuals sum to zero on each block. Since \(h(m_{n,\mathrm{iso}}^\star(X_i))=h(c_j)\) for all \(i \in B_j\),
\[
\sum_{i=1}^n h(m_{n,\mathrm{iso}}^\star(X_i))\{Y_i-m_{n,\mathrm{iso}}^\star(X_i)\}
=
\sum_{j=1}^J h(c_j)\sum_{i \in B_j}(Y_i-c_j)
=
0.
\]
\end{proof}

\begin{proof}[Proof of \cref{prop:iso-repr}]
This directly follows from Lemma~\ref{lem:block}.
\end{proof}

\begin{proof}[Proof of \cref{prop:linear}]
The normal equations for \cref{eq:linear-cal} are
\[
\Pl\{Y-m_{n,\mathrm{lin}}^\star(X)\}=0
\qquad\text{and}\qquad
\Pl\bigl[\mhat(X)\{Y-m_{n,\mathrm{lin}}^\star(X)\}\bigr]=0.
\]
Because \(m_{n,\mathrm{lin}}^\star(X)=\hat a_{\mathrm{lin}}\mhat(X)+\hat b_{\mathrm{lin}}\), these imply
\[
\Pl\bigl[m_{n,\mathrm{lin}}^\star(X)\{Y-m_{n,\mathrm{lin}}^\star(X)\}\bigr]=0.
\]
Hence, for any \(\hat w(X)=a+b\,m_{n,\mathrm{lin}}^\star(X)\),
\[
\Pl\bigl[\hat w(X)\{Y-m_{n,\mathrm{lin}}^\star(X)\}\bigr]
=
a\,\Pl\{Y-m_{n,\mathrm{lin}}^\star(X)\}
+
b\,\Pl\bigl[m_{n,\mathrm{lin}}^\star(X)\{Y-m_{n,\mathrm{lin}}^\star(X)\}\bigr]
=
0.
\]
Adding this zero term to \cref{eq:linear-plugin} yields \cref{eq:linear-ppi}.
\end{proof}

\section{Proofs of the efficient influence pairs}
\label{app:eif}

\begin{proof}[Proof of \cref{prop:eif-full}]
Under the joint two-sample experiment \(\mathbb D^{(n,N)}(P)=P^{\otimes n}\otimes P_X^{\otimes N}\), the tangent space is the product of the tangent space for the labeled sample and the tangent space for the unlabeled sample, with inner product weighted by the limiting sample fractions \(\rhozero\) and \(1-\rhozero\). Perturbations of the marginal law of \(X\) affect both samples, while perturbations of the conditional law of \(Y \mid X\) appear only through the labeled sample. The pathwise derivative of the parameter \(\Psi\{\mathbb D^{(n,N)}(P)\}=\E_P[Y]=\E_{P_X}[\muzero(X)]\) is therefore represented by the pair
\[
\muzero(X)-\psi(P)+\rhozero^{-1}\{Y-\muzero(X)\}
\quad\text{and}\quad
\muzero(\widetilde X)-\psi(P),
\]
which yields \cref{eq:eif-full-L,eq:eif-full-U}.
\end{proof}

\begin{proof}[Proof of \cref{prop:eif-red}]
Under the reduced experiment \(\mathcal E_{m_0}\), the observed data are \((m_0(X),Y)\) in the labeled sample and \(m_0(\widetilde X)\) in the unlabeled sample. On the coarsened submodel defined by the identifying restriction \(\E_P[Y]=\E_P[m_0(X)]\), the target \(\psi(P)=\E_P[Y]\) is equivalently the mean of the observed summary \(m_0(X)\). The standard two-sample efficient influence pair for this semisupervised mean functional is therefore
\[
m_0(X)-\psi_0+\rhozero^{-1}\{Y-m_0(X)\}
\quad\text{and}\quad
m_0(\widetilde X)-\psi_0,
\]
which gives \cref{eq:eif-red-L,eq:eif-red-U}.
\end{proof}

\begin{lemma}[Empirical calibration centers the limit]
\label{lem:center}
Suppose \cref{assump:design,assump:isoreg} hold and the calibration score equations contain the constant score \(h\equiv 1\). Then
\[
\E_{\Px}[m_0(X)]=\psi_0.
\]
\end{lemma}

\textbf{Proof convention.} In the following proofs, fix \(\delta>0\). By \(\sup_x |m_{n,\mathrm{iso}}^\star(x)|=O_P(1)\), there exist a deterministic constant \(C=C_\delta\ge C_0\) and \(n_0<\infty\) such that the event
\[
A_{n,C}:=\Bigl\{\sup_x |m_{n,\mathrm{iso}}^\star(x)|\le C\Bigr\}
\]
satisfies \(P_0(A_{n,C})\ge 1-\delta\) for all \(n\ge n_0\). We carry out the deterministic bounded-class arguments on \(A_{n,C}\), where all implicit constants may depend on \(C\) but not on \(n\). Since \(\delta\) is arbitrary, the resulting rates and distributional statements hold unconditionally.


\begin{lemma}[Isotonic \(L^2\) rate]
\label{lem:isorate}
Suppose \cref{assump:design,assump:isoreg} holds. Then
\[
\|m_{n,\mathrm{iso}}^\star-m_0\|_{2,\Px}^2=O_p(n^{-2/3}).
\]
\end{lemma}



\begin{proof}[Proof of \cref{lem:isorate}]
 


Let \(T:=\mhat(X)\) and let
\[
\mathcal M_C:=\{f(T):f\in\mathcal F_{\mathrm{iso}},\ \|f(T)\|_\infty\le C\}.
\]
By \cref{assump:isoreg}, \(g_0(T)\) is bounded by \(C_0\le C\), and on \(A_{n,C}\) we have \(m_{n,\mathrm{iso}}^\star\in\mathcal M_C\). Since clipping any monotone candidate to \([-C,C]\) preserves monotonicity and can only decrease squared risk against \(g_0(T)\), the population isotonic projection \(m_0\) also belongs to \(\mathcal M_C\). Thus the problem is a one-dimensional isotonic least-squares regression of \(Y\) on the score \(T\). Let
\[
g_0(T):=\E[Y\mid T].
\]
Then, for every \(m\in\mathcal M_C\),
\[
P_0\{(Y-m)^2\}
=
P_0\{(Y-g_0)^2\}+P_0\{(g_0-m)^2\},
\]
which shows that \(m_0\) is the \(L^2(P_T)\)-projection of \(g_0\) onto the closed convex set \(\mathcal M_C\).

Now define the centered class
\[
\mathcal F_0:=\{g=m-m_0:\ m\in\mathcal M_C\}.
\]
For \(g\in\mathcal F_0\), let \(\ell_g(Z):=(Y-m_0(T)-g(T))^2-(Y-m_0(T))^2\), where \(Z=(X,Y)\). Since \(m_{n,\mathrm{iso}}^\star\) minimizes the empirical squared risk over \(\mathcal F_{\mathrm{iso}}\), on \(A_{n,C}\) the random element \(\hat g:=m_{n,\mathrm{iso}}^\star-m_0\) belongs to \(\mathcal F_0\) and satisfies
\[
\Pl \ell_{\hat g}\le 0=\Pl \ell_0.
\]
Moreover,
\[
P_0 \ell_g
=
\|g\|_{2,\Px}^2-2P_0\left[(Y-m_0)g\right].
\]
Because \(g\) is measurable with respect to \(T\), the tower property gives
\[
P_0\left[(Y-m_0)g\right]
=
P_0\left[(g_0-m_0)g\right].
\]
Since \(m_0\) is the \(L^2(P_T)\)-projection of \(g_0\) onto the convex set \(\mathcal M_C\), the Hilbert-space projection inequality yields
\[
P_0\left[(g_0-m_0)(m-m_0)\right]\le 0
\qquad\text{for every }m\in\mathcal M_C.
\]
Hence
\[
P_0 \ell_g \ge \|g\|_{2,\Px}^2,
\]
so squared loss has the required quadratic margin around \(m_0\).

We next bound the modulus of continuity of the empirical process indexed by the localized classes
\[
\mathcal F_\delta:=\{g\in\mathcal F_0:\|g\|_{2,\Px}\le \delta\}.
\]
Because every \(m\in\mathcal M_C\) is bounded by \(C\), every \(g\in\mathcal F_\delta\) is uniformly bounded, and the standard bracketing bound for one-dimensional monotone classes gives
\[
\log N_{[]}(\varepsilon,\mathcal F_\delta,L_2(\Px))
\lesssim
\delta/\varepsilon,
\qquad 0<\varepsilon<\delta.
\]
Therefore
\[
J_{[]}(\delta,\mathcal F_\delta,L_2(\Px))\lesssim \delta^{1/2};
\]
see, e.g., \citet{vdvwellner1996}.

Write \(\xi:=Y-m_0(T)\), so that \(\ell_g=g^2-2\xi g\). The class \(\{g^2:g\in\mathcal F_\delta\}\) is a bounded Lipschitz image of \(\mathcal F_\delta\), so by Lemma 3.4.2 of \citet{vdvwellner1996} and the same entropy bound,
\[
E^*\sup_{g\in\mathcal F_\delta}\sqrt{n}\,|(\Pl-P_0)(g^2)|
\lesssim
\delta^{1/2}.
\]
For the multiplier term, \cref{assump:isoreg} implies that \(Y-g_0(T)\) is sub-Gaussian or subexponential. Since both \(g_0(T)\) and \(m_0(T)\) are bounded by \(C\), the difference \(g_0(T)-m_0(T)\) is bounded, and therefore \(\xi=Y-m_0(T)\) is also sub-Gaussian or subexponential. Thus the product class \(\xi\mathcal F_\delta:=\{\xi g:g\in\mathcal F_\delta\}\) satisfies the same localized bracketing bound in the Bernstein-type norm used in the subexponential least-squares proof of \citet{Bibaut2019}, yielding
\[
E^*\sup_{g\in\mathcal F_\delta}\sqrt{n}\,|(\Pl-P_0)(\xi g)|
\lesssim
\delta^{1/2}.
\]
Combining the previous two displays,
\[
E^*\sup_{g\in\mathcal F_\delta}\sqrt{n}\,|(\Pl-P_0)\ell_g|
\lesssim
\delta^{1/2}.
\]
We may therefore apply Theorem 3.4.1 of \citet{vdvwellner1996} with \(d(g)=\|g\|_{2,\Px}\), quadratic margin \(P_0\ell_g\gtrsim d^2(g)\), and modulus \(\phi_n(\delta)\asymp\delta^{1/2}\). The fixed-point condition \(r_n^2\phi_n(r_n^{-1})\lesssim \sqrt n\) gives \(r_n\asymp n^{1/3}\), and hence
\[
\|m_{n,\mathrm{iso}}^\star-m_0\|_{2,\Px}=O_p(n^{-1/3}),
\qquad
\|m_{n,\mathrm{iso}}^\star-m_0\|_{2,\Px}^2=O_p(n^{-2/3}).
\]
The preceding display is obtained on \(A_{n,C}\). Since \(P_0(A_{n,C}^c)\le \delta\) for all large \(n\), and \(\delta>0\) was arbitrary, the same rate holds unconditionally. This is the classical one-dimensional isotonic rate, stated directly in the random-design \(L^2(\Px)\) norm relevant for the theorem.
\end{proof}

\begin{proof}[Proof of \cref{lem:center}]
Because \(h\equiv 1\) belongs to the calibration score class, \cref{eq:cal-score} gives
\[
\Pl\{Y-m_{n,\mathrm{iso}}^\star(X)\}=0.
\]
Hence
\[
\Pl Y-\Pl m_0(X)=\Pl\{m_{n,\mathrm{iso}}^\star(X)-m_0(X)\}=o_p(1),
\]
where the last step follows from \cref{lem:isorate} and Cauchy--Schwarz. By \cref{assump:design}, the law of large numbers yields \(\Pl Y\to_p \psi_0\) and \(\Pl m_0(X)\to_p \E_{\Px}[m_0(X)]\). Therefore \(\E_{\Px}[m_0(X)]=\psi_0\).
\end{proof}

\begin{lemma}[Centered empirical-process bounds for isotonic calibration]
\label{lem:eprate}
Suppose \cref{assump:design,assump:isoreg} holds. Then
\[
(\Pl-\Px)\{m_{n,\mathrm{iso}}^\star(X)-m_0(X)\}=O_p(n^{-2/3}),
\qquad
(\Pu-\Px)\{m_{n,\mathrm{iso}}^\star(\widetilde X)-m_0(\widetilde X)\}=O_p\bigl(n^{-1/3}N^{-1/2}\bigr).
\]
\end{lemma}

\begin{proof}[Proof of \cref{lem:eprate}]
By \cref{lem:isorate}, the isotonic fit satisfies
\[
\|m_{n,\mathrm{iso}}^\star-m_0\|_{2,\Px}^2=O_p(n^{-2/3}).
\]
To control the centered empirical-process terms, consider the localized class
\[
\mathcal F_{n,C}:=\{g=f\circ\mhat-m_0:\ f\in\mathcal F_{\mathrm{iso}},\ \|f\circ\mhat\|_\infty\le C,\ \|g\|_{2,\Px}^2\lesssim n^{-2/3}\}.
\]
On \(A_{n,C}\), the function \(m_{n,\mathrm{iso}}^\star-m_0\) belongs to \(\mathcal F_{n,C}\). Moreover, \(\mathcal F_{n,C}\) has a bounded envelope because both \(f\circ\mhat\) and \(m_0\) are bounded by \(C\). Because \(\mhat\) is fixed, \(\mathcal F_{n,C}\) is a localized class of bounded monotone transformations of a one-dimensional score, and standard entropy bounds for monotone classes imply a finite entropy integral. Applying a local maximal inequality for empirical processes, such as Lemma 3.4.2 of \citet{vdvwellner1996}, at localization radius \(n^{-1/3}\) then gives, on \(A_{n,C}\),
\[
|(\Pl-\Px)\{m_{n,\mathrm{iso}}^\star-m_0\}|=O_p(n^{-2/3}).
\]
Since \(P_0(A_{n,C}^c)\le \delta\) for all large \(n\), and \(\delta>0\) was arbitrary, the same bound holds unconditionally.
For the unlabeled sample, conditional on the labeled data, \(m_{n,\mathrm{iso}}^\star-m_0\) is fixed and the unlabeled sample is independent, so
\[
(\Pu-\Px)\{m_{n,\mathrm{iso}}^\star-m_0\}
=
O_p\left(N^{-1/2}\|m_{n,\mathrm{iso}}^\star-m_0\|_{2,\Px}\right)
=
O_p\bigl(n^{-1/3}N^{-1/2}\bigr)
\]
by \cref{lem:isorate}.
\end{proof}

\section{Proof of asymptotic linearity}
\label{app:al}

\begin{proof}[Proof of \cref{thm:al}]
By \cref{lem:center}, \(\Px m_0=\psi_0\). Using \cref{thm:repr} with \(\what \equiv 1\),
\[
\widehat{\psi}_{\mathrm{iso}}
=
\rhon \Pl\{m_{n,\mathrm{iso}}^\star(X)\}+(1-\rhon)\Pu\{m_{n,\mathrm{iso}}^\star(\widetilde X)\}+\rhon \Pl\{Y-m_{n,\mathrm{iso}}^\star(X)\}.
\]
Add and subtract the AIPW form built from \(m_0\):
\[
\widehat{\psi}_{\mathrm{iso}}
=
\rhon \Pl\left\{m_0+\rhon^{-1}(Y-m_0)\right\}
+
(1-\rhon)\Pu\{m_0(\widetilde X)\}
+
(1-\rhon)(\Pu-\Pl)\{m_{n,\mathrm{iso}}^\star-m_0\}.
\]
Subtracting \(\psi_0\) gives
\[
\begin{aligned}
\widehat{\psi}_{\mathrm{iso}}-\psi_0
&=
\rhon(\Pl-P_0)\Bigl\{m_0-\psi_0+\rhon^{-1}(Y-m_0)\Bigr\}\\
&\quad +
(1-\rhon)(\Pu-\Px)(m_0-\psi_0)\\
&\quad +
(1-\rhon)(\Pu-\Pl)\{m_{n,\mathrm{iso}}^\star-m_0\}.
\end{aligned}
\]
To replace \(\rhon^{-1}\) by \(\rhozero^{-1}\), note that
\[
\rhon(\Pl-P_0)\Bigl\{(\rhon^{-1}-\rhozero^{-1})(Y-m_0)\Bigr\}
=
\Bigl(1-\frac{\rhon}{\rhozero}\Bigr)(\Pl-P_0)(Y-m_0).
\]
The right-hand side is \(o_p(M^{-1/2})\) because \(\rhon\to\rhozero\) and \((\Pl-P_0)(Y-m_0)=O_p(n^{-1/2})=O_p(M^{-1/2})\). Also,
\[
(1-\rhon)(\Pu-\Pl)\{m_{n,\mathrm{iso}}^\star-m_0\}
=(1-\rhon)(\Pu-\Px)\{m_{n,\mathrm{iso}}^\star-m_0\}-(1-\rhon)(\Pl-\Px)\{m_{n,\mathrm{iso}}^\star-m_0\},
\]
so the labeled term is \(O_p(n^{-2/3})\) and the unlabeled term is \(O_p(n^{-1/3}N^{-1/2})\) by \cref{lem:eprate}. Thus the remainder \(R_{n,N}\), defined as the difference between \(\widehat{\psi}_{\mathrm{iso}}-\psi_0\) and the two leading empirical-process terms in \cref{eq:al}, satisfies the displayed bound in \cref{thm:al}. Since each component is \(o_p(M^{-1/2})\), this also yields \cref{eq:al}.
\end{proof}

\begin{proof}[Proof of \cref{cor:normal}]
By \cref{thm:al},
\[
\sqrt{M}\,(\widehat{\psi}_{\mathrm{iso}}-\psi_0)
=
\frac{1}{\sqrt{M}}\sum_{i=1}^n D_{m_0}^L(X_i,Y_i)
+
\frac{1}{\sqrt{M}}\sum_{j=1}^N D_{m_0}^U(\widetilde X_j)
+
o_p(1).
\]
The labeled and unlabeled samples are independent, each summand has mean zero, and \(\rhon\to\rhozero\in(0,1)\) by \cref{assump:design}. A triangular-array central limit theorem therefore gives the stated asymptotic normal law with variance \(\sigma_0^2\).

For the variance estimator, \cref{lem:isorate} implies \(\Pl\{(m_{n,\mathrm{iso}}^\star-m_0)^2\}=o_p(1)\), and \(\widehat{\psi}_{\mathrm{iso}}\to_p\psi_0\) follows from \cref{thm:al}. Hence
\[
\widehat D_i^L-D_{m_0}^L(X_i,Y_i)=o_p(1)
\qquad\text{and}\qquad
\widehat D_j^U-D_{m_0}^U(\widetilde X_j)=o_p(1)
\]
in empirical \(L^2\), which yields \(\widehat\sigma^2\to_p \sigma_0^2\). The Wald interval claim then follows by Slutsky's theorem.
\end{proof}

\section{Proof of reduced-model efficiency}
\label{app:eff}

\begin{proof}[Proof of \cref{cor:efficiency}]
Part (i) follows because, under the conditions of \cref{thm:al}, the influence pair in \cref{thm:al} is exactly the pair in \cref{prop:eif-red}. Part (ii) follows because if \(m_0=\muzero\) almost surely, then the same influence pair also coincides with the efficient influence pair in \cref{prop:eif-full}. Part (iii) follows because monotonicity of \(\E[Y\mid m(X)]\) in \(m(X)\) identifies the score-based regression with the isotonic target \(m_0\), so the reduced model indexed by \(m_0\) is the same as the model indexed by the original score \(m(X)\).
\end{proof}
 

\section{Pooled i.i.d.\ Missing-Data Formulation}
\label{app:iid}

The joint two-sample setup can be embedded into a pooled i.i.d.\ model by introducing a sample-membership indicator \(R \in \{0,1\}\), where \(R=1\) denotes labeled observations and \(R=0\) denotes unlabeled observations. In that formulation one observes i.i.d.\ data \(O=(X,R,RY)\) with \(\Prob(R=1)=\rhozero\), and the calibrated estimator becomes
\[
\psihat = \mathbb P_M(\mtilde) + \mathbb P_M\left[\frac{R}{\rhozero}\{Y-\mtilde\}\right]
\]
up to the same calibration identity as in the main text. The corresponding observed-data efficient influence function is
\[
m_0(X)-\psi_0+\frac{R}{\rhozero}\{Y-m_0(X)\},
\]
which is the single-sample analogue of the reduced-model influence pair in \cref{prop:eif-red}. Thus the usual i.i.d.\ missing-data formulation is a convenient special case, but the main text keeps the two-dataset structure explicit.

 

\section{Additional simulation details}
\label{app:simulation-extra}
The mildly miscalibrated synthetic regime serves mainly as a sanity check: once the original score is already close to calibrated, all score-based estimators are nearly indistinguishable, so there is little to gain from extra flexibility. In that regime, AIPW and PPI also track each other closely, with the clearest separation again appearing when the labeled sample is a non-negligible part of the pooled sample. The main role of the synthetic study is therefore to show that stronger nonlinear miscalibration is exactly the regime in which recalibration becomes useful, not to suggest that recalibration should be applied indiscriminately.

\begin{figure}[h]
\centering
\includegraphics[width=\textwidth]{exp_plots_tables/sim_fig_bias_coverage.pdf}
\caption{Additional synthetic diagnostics at the largest unlabeled ratio \(N/n\). The top row corresponds to the mildly miscalibrated regime and shows that the score-based estimators nearly overlap. The bottom row echoes the main-text advantage of recalibration under stronger nonlinear distortion.}
\label{fig:simulation-appendix}
\end{figure}

\section{Additional empirical details}
\label{app:empirical}
This appendix records the benchmark construction used in \cref{sec:empirical}. We use the official \texttt{ppi\_py} real-data mean-estimation examples and preserve their labeled-sample-size grids. The \texttt{forest} experiment uses \(n \in \{50,100,\dots,500\}\), while \texttt{galaxies} and semisupervised \texttt{census\_income} use \(n \in \{50,155,261,366,472,577,683,788,894,1000\}\). For each \(n\), the remaining observations form the unlabeled sample. We rerandomize this split 5000 times at each \(n\), always evaluating all estimators on the same split and taking the full-sample mean as the benchmark truth.

The estimator set mirrors the original PPI comparison wherever possible. We report the labeled-only estimator, the imputation benchmark on binary-outcome tasks, the official semisupervised linear baseline from \texttt{ppi\_py},  PPI computed using the official \texttt{ppi\_py} implementation, and the classical AIPW estimator based on the same prediction score. We then add four calibration-based plug-in estimators: affine linear recalibration; Platt scaling on the binary-outcome benchmarks; isotonic recalibration with a fixed minimum bin size of 10; and Venn--Abers shrinkage on the binary-outcome benchmarks. Platt scaling and Venn--Abers are omitted on \texttt{census\_income} because that outcome is not binary. As discussed in the main text, PPI++ and prognostic-score adjustment are not shown as separate baselines because they coincide algebraically with linear recalibration in this one-dimensional score setting.

For the calibration-based plug-in estimators, we fit the recalibrator on the labeled pairs \(\{(\mhat(X_i),Y_i)\}_{i=1}^n\), transform both the labeled and unlabeled scores, and then compute the pooled plug-in mean together with the same semisupervised influence-style Wald interval used throughout the benchmark code. AIPW uses the same score but averages it over the pooled covariate sample before adding the labeled residual correction, whereas  PPI uses only the unlabeled score average. In theory this retains less labeled-sample score information, though in the finite-sample benchmark results the empirical variance gap is often small and not uniform across all datasets. Linear recalibration regresses \(Y\) on an intercept and the score, then clips the fitted values to the range of predicted outcomes observed in the labeled data to prevent extrapolation on the unlabeled data. Platt scaling fits a logistic regression of the binary outcome on a stabilized logit transform of the score. The fixed-bin isotonic estimator also clips to the labeled outcome range. The reported summaries are Monte Carlo averages of bias, empirical variance, MSE, interval coverage, and relative efficiency versus PPI across the repeated random splits. The complete numerical outputs, including dataset-level tables and per-metric plots, are generated automatically by the accompanying reproduction pipeline.

\begin{table}[h]
\centering
\small
\resizebox{\textwidth}{!}{\begin{tabular}{lllrccccc}
\toprule
Dataset & $n$ & $N$ & Estimator & Bias & Variance & MSE & Coverage & RelEff \\
\midrule
census\_income\_semisupervised\_mean & 50 & 380041 & AIPW & 64.1415 & 74530505.1597 & 74534619.2853 & 0.802 & 1.000 \\
census\_income\_semisupervised\_mean & 50 & 380041 & Classical & 174.3296 & 112879690.3397 & 112910081.1357 & 0.841 & 0.660 \\
census\_income\_semisupervised\_mean & 50 & 380041 & IsoCal (min bin 10) & -479.6340 & 75379435.8991 & 75609484.6737 & 0.777 & 0.989 \\
census\_income\_semisupervised\_mean & 50 & 380041 & LinearCal & -22.9713 & 72975497.5814 & 72976025.2610 & 0.784 & 1.021 \\
census\_income\_semisupervised\_mean & 50 & 380041 & PPI & 64.1270 & 74528348.5082 & 74532460.7743 & 0.802 & 1.000 \\
census\_income\_semisupervised\_mean & 50 & 380041 & Semisup. linear & 49.6769 & 106851003.3081 & 106853471.1059 & 0.844 & 0.697 \\
census\_income\_semisupervised\_mean & 577 & 379514 & AIPW & 1.6631 & 6451762.8714 & 6451765.6373 & 0.890 & 1.000 \\
census\_income\_semisupervised\_mean & 577 & 379514 & Classical & -44.6870 & 9708740.7725 & 9710737.7047 & 0.894 & 0.664 \\
census\_income\_semisupervised\_mean & 577 & 379514 & IsoCal (min bin 10) & -23.7251 & 6409823.1141 & 6410385.9947 & 0.880 & 1.006 \\
census\_income\_semisupervised\_mean & 577 & 379514 & LinearCal & 13.6244 & 6307041.6066 & 6307227.2310 & 0.883 & 1.023 \\
census\_income\_semisupervised\_mean & 577 & 379514 & PPI & 1.7336 & 6449769.3295 & 6449772.3348 & 0.890 & 1.000 \\
census\_income\_semisupervised\_mean & 577 & 379514 & Semisup. linear & -19.0374 & 8948168.5056 & 8948530.9279 & 0.894 & 0.721 \\
census\_income\_semisupervised\_mean & 1000 & 379091 & AIPW & -20.5346 & 3817516.0488 & 3817937.7195 & 0.891 & 0.999 \\
census\_income\_semisupervised\_mean & 1000 & 379091 & Classical & -11.5142 & 5706667.9025 & 5706800.4785 & 0.898 & 0.669 \\
census\_income\_semisupervised\_mean & 1000 & 379091 & IsoCal (min bin 10) & -40.5935 & 3746139.4177 & 3747787.2512 & 0.879 & 1.018 \\
census\_income\_semisupervised\_mean & 1000 & 379091 & LinearCal & -27.7615 & 3679849.2001 & 3680619.8996 & 0.890 & 1.037 \\
census\_income\_semisupervised\_mean & 1000 & 379091 & PPI & -20.5584 & 3815435.7633 & 3815858.4117 & 0.891 & 1.000 \\
census\_income\_semisupervised\_mean & 1000 & 379091 & Semisup. linear & -17.5832 & 5259398.9756 & 5259708.1447 & 0.898 & 0.725 \\
forest\_mean & 50 & 1546 & AIPW & -0.0007 & 0.0018 & 0.0018 & 0.934 & 1.003 \\
forest\_mean & 50 & 1546 & Classical & 0.0002 & 0.0025 & 0.0025 & 0.948 & 0.709 \\
forest\_mean & 50 & 1546 & Imputation & -0.0733 & 0.0000 & 0.0054 & 0.000 & inf \\
forest\_mean & 50 & 1546 & IsoCal (min bin 10) & -0.0035 & 0.0020 & 0.0020 & 0.919 & 0.895 \\
forest\_mean & 50 & 1546 & LinearCal & -0.0012 & 0.0018 & 0.0018 & 0.922 & 0.957 \\
forest\_mean & 50 & 1546 & Platt & -0.0002 & 0.0019 & 0.0019 & 0.915 & 0.946 \\
forest\_mean & 50 & 1546 & PPI & -0.0007 & 0.0018 & 0.0018 & 0.935 & 1.000 \\
forest\_mean & 50 & 1546 & Venn-Abers & -0.0026 & 0.0019 & 0.0019 & 0.899 & 0.909 \\
forest\_mean & 300 & 1296 & AIPW & -0.0001 & 0.0002 & 0.0002 & 0.974 & 1.042 \\
forest\_mean & 300 & 1296 & Classical & -0.0003 & 0.0004 & 0.0004 & 0.967 & 0.723 \\
forest\_mean & 300 & 1296 & Imputation & -0.0733 & 0.0000 & 0.0054 & 0.000 & inf \\
forest\_mean & 300 & 1296 & IsoCal (min bin 10) & -0.0002 & 0.0003 & 0.0003 & 0.966 & 1.024 \\
forest\_mean & 300 & 1296 & LinearCal & -0.0002 & 0.0002 & 0.0002 & 0.971 & 1.030 \\
forest\_mean & 300 & 1296 & Platt & -0.0001 & 0.0002 & 0.0002 & 0.971 & 1.033 \\
forest\_mean & 300 & 1296 & PPI & -0.0001 & 0.0003 & 0.0003 & 0.973 & 1.000 \\
forest\_mean & 300 & 1296 & Venn-Abers & -0.0002 & 0.0003 & 0.0003 & 0.966 & 1.024 \\
forest\_mean & 500 & 1096 & AIPW & 0.0000 & 0.0001 & 0.0001 & 0.988 & 1.169 \\
forest\_mean & 500 & 1096 & Classical & 0.0003 & 0.0002 & 0.0002 & 0.984 & 0.849 \\
forest\_mean & 500 & 1096 & Imputation & -0.0733 & 0.0000 & 0.0054 & 0.000 & inf \\
forest\_mean & 500 & 1096 & IsoCal (min bin 10) & 0.0001 & 0.0001 & 0.0001 & 0.987 & 1.171 \\
forest\_mean & 500 & 1096 & LinearCal & 0.0000 & 0.0001 & 0.0001 & 0.988 & 1.173 \\
forest\_mean & 500 & 1096 & Platt & 0.0000 & 0.0001 & 0.0001 & 0.988 & 1.177 \\
forest\_mean & 500 & 1096 & PPI & -0.0001 & 0.0001 & 0.0001 & 0.985 & 1.000 \\
forest\_mean & 500 & 1096 & Venn-Abers & 0.0001 & 0.0001 & 0.0001 & 0.986 & 1.177 \\
galaxies\_mean & 50 & 16693 & AIPW & 0.0003 & 0.0018 & 0.0018 & 0.888 & 1.000 \\
galaxies\_mean & 50 & 16693 & Classical & 0.0004 & 0.0040 & 0.0040 & 0.883 & 0.449 \\
galaxies\_mean & 50 & 16693 & Imputation & -0.0433 & 0.0000 & 0.0019 & 0.000 & inf \\
galaxies\_mean & 50 & 16693 & IsoCal (min bin 10) & -0.0040 & 0.0020 & 0.0020 & 0.847 & 0.897 \\
galaxies\_mean & 50 & 16693 & LinearCal & -0.0001 & 0.0019 & 0.0019 & 0.872 & 0.970 \\
galaxies\_mean & 50 & 16693 & Platt & 0.0002 & 0.0019 & 0.0019 & 0.858 & 0.965 \\
galaxies\_mean & 50 & 16693 & PPI & 0.0003 & 0.0018 & 0.0018 & 0.888 & 1.000 \\
galaxies\_mean & 50 & 16693 & Venn-Abers & -0.0001 & 0.0019 & 0.0019 & 0.832 & 0.933 \\
galaxies\_mean & 577 & 16166 & AIPW & -0.0002 & 0.0001 & 0.0001 & 0.908 & 1.003 \\
galaxies\_mean & 577 & 16166 & Classical & -0.0000 & 0.0003 & 0.0003 & 0.910 & 0.456 \\
galaxies\_mean & 577 & 16166 & Imputation & -0.0433 & 0.0000 & 0.0019 & 0.000 & inf \\
galaxies\_mean & 577 & 16166 & IsoCal (min bin 10) & -0.0002 & 0.0001 & 0.0001 & 0.900 & 0.989 \\
galaxies\_mean & 577 & 16166 & LinearCal & -0.0002 & 0.0001 & 0.0001 & 0.906 & 1.001 \\
galaxies\_mean & 577 & 16166 & Platt & -0.0003 & 0.0001 & 0.0001 & 0.905 & 0.999 \\
galaxies\_mean & 577 & 16166 & PPI & -0.0003 & 0.0001 & 0.0001 & 0.908 & 1.000 \\
galaxies\_mean & 577 & 16166 & Venn-Abers & -0.0002 & 0.0001 & 0.0001 & 0.899 & 0.991 \\
galaxies\_mean & 1000 & 15743 & AIPW & -0.0002 & 0.0001 & 0.0001 & 0.913 & 1.004 \\
galaxies\_mean & 1000 & 15743 & Classical & -0.0001 & 0.0002 & 0.0002 & 0.899 & 0.444 \\
galaxies\_mean & 1000 & 15743 & Imputation & -0.0433 & 0.0000 & 0.0019 & 0.000 & inf \\
galaxies\_mean & 1000 & 15743 & IsoCal (min bin 10) & -0.0002 & 0.0001 & 0.0001 & 0.907 & 0.987 \\
galaxies\_mean & 1000 & 15743 & LinearCal & -0.0002 & 0.0001 & 0.0001 & 0.913 & 1.000 \\
galaxies\_mean & 1000 & 15743 & Platt & -0.0002 & 0.0001 & 0.0001 & 0.910 & 0.996 \\
galaxies\_mean & 1000 & 15743 & PPI & -0.0002 & 0.0001 & 0.0001 & 0.912 & 1.000 \\
galaxies\_mean & 1000 & 15743 & Venn-Abers & -0.0002 & 0.0001 & 0.0001 & 0.905 & 0.984 \\
\bottomrule
\end{tabular}
}
\caption{Representative numerical summary at the smallest, middle, and largest labeled sample sizes for each reproduced benchmark. The reported quantities are Monte Carlo bias, empirical variance of the point estimator, MSE, Wald-interval coverage, and relative efficiency versus PPI.}
\label{tab:ppi-benchmarks-full}
\end{table}

\begin{figure}[h]
\centering
\includegraphics[width=\textwidth]{exp_plots_tables/fig_paper_metric_grid.pdf}
\caption{Full diagnostic grid for the reproduced PPI benchmarks, showing bias, empirical variance, MSE, coverage, and relative efficiency versus PPI across labeled-sample-size regimes. The dashed horizontal line in the coverage panels marks the nominal target coverage \(1-\alpha\), and the dashed horizontal line in the relative-efficiency panels marks parity with  PPI. Platt scaling and Venn--Abers appear only on the binary-outcome datasets where they are well defined.}
\label{fig:ppi-benchmarks-full}
\end{figure}





\section{\texttt{Python} Code}
\label{appendix:code}

\subsection{Isotonic regression via XGBoost}
\begin{codeblock}
import numpy as np
import xgboost as xgb

def isoreg_with_xgboost(x, y, max_depth=15, min_child_weight=20, weights=None):
    """
    Fit a monotone calibrator f so that f(x) is nondecreasing in x.

    Parameters
    ----------
    x : array-like, shape (n,) or (n, 1)
        Predictor used for calibration. In the PPI application this is usually
        the score m(X), so x is typically one-dimensional.
    y : array-like, shape (n,)
        Labeled outcomes.
    max_depth : int, default=15
        Maximum depth for the one-round XGBoost fit.
    min_child_weight : float, default=20
        Minimum total weight in a terminal node.
    weights : array-like, optional
        Optional observation weights.

    Returns
    -------
    predict_fn : callable
        Function mapping new x values to calibrated predictions f(x).
    """
    x = np.asarray(x)
    y = np.asarray(y).reshape(-1)
    x = x.reshape(len(y), -1)

    data = xgb.DMatrix(data=x, label=y, weight=weights)
    params = {
        "max_depth": max_depth,
        "min_child_weight": min_child_weight,
        "monotone_constraints": "(" + ",".join(["1"] * x.shape[1]) + ")",
        "eta": 1.0,
        "gamma": 0.0,
        "lambda": 0.0,
        "objective": "reg:squarederror",
        "verbosity": 0,
    }
    iso_fit = xgb.train(params=params, dtrain=data, num_boost_round=1)

    def predict_fn(x_new):
        x_new = np.asarray(x_new)
        if x_new.ndim == 1:
            x_new = x_new.reshape(-1, 1)
        data_pred = xgb.DMatrix(data=x_new)
        return iso_fit.predict(data_pred)

    return predict_fn
\end{codeblock}

\subsection{Isotonic-calibrated plug-in}
\begin{codeblock}
import numpy as np

from statistics import NormalDist


def _as_1d_array(x, name):
    x = np.asarray(x).reshape(-1)
    if x.size == 0:
        raise ValueError(f"{name} must be non-empty.")
    return x


def _resolve_scores(
    X_labeled=None,
    Y_labeled=None,
    X_unlabeled=None,
    score_fn=None,
    score_labeled=None,
    score_unlabeled=None,
):
    """
    Resolve the initial score m(X) on the labeled and unlabeled samples.

    Users may either supply:
    1. a scoring function `score_fn`, together with `X_labeled` and
       `X_unlabeled`, or
    2. precomputed scores `score_labeled` and `score_unlabeled`.
    """
    if score_fn is not None:
        if X_labeled is None or X_unlabeled is None:
            raise ValueError(
                "If score_fn is provided, then X_labeled and X_unlabeled "
                "must also be provided."
            )
        score_labeled = score_fn(X_labeled)
        score_unlabeled = score_fn(X_unlabeled)
    elif score_labeled is None or score_unlabeled is None:
        raise ValueError(
            "Provide either score_fn with X_labeled/X_unlabeled, or provide "
            "both score_labeled and score_unlabeled."
        )

    score_labeled = _as_1d_array(score_labeled, "score_labeled")
    score_unlabeled = _as_1d_array(score_unlabeled, "score_unlabeled")

    if Y_labeled is not None and len(score_labeled) != len(Y_labeled):
        raise ValueError("score_labeled and Y_labeled must have the same length.")

    return score_labeled, score_unlabeled


def calibrated_plugin_ppi(
    Y_labeled,
    X_labeled=None,
    X_unlabeled=None,
    score_fn=None,
    score_labeled=None,
    score_unlabeled=None,
    min_obs=20,
    alpha=0.05,
    inference="wald",
    n_bootstrap=999,
    random_state=None,
    max_depth=15,
):
    """
    Calibrated plug-in PPI estimator for psi = E[Y].

    The estimator first calibrates the initial score m(X) on the labeled data,
    then averages the calibrated score over the pooled labeled and unlabeled
    covariate sample.

    Parameters
    ----------
    Y_labeled : array-like, shape (n,)
        Outcomes in the labeled sample.
    X_labeled : array-like, optional
        Covariates in the labeled sample. Needed only when score_fn is used.
    X_unlabeled : array-like, optional
        Covariates in the unlabeled sample. Needed only when score_fn is used.
    score_fn : callable, optional
        Function returning the initial score m(X).
    score_labeled : array-like, optional
        Precomputed values m(X) on the labeled sample.
    score_unlabeled : array-like, optional
        Precomputed values m(Xtilde) on the unlabeled sample.
    min_obs : int, default=20
        Minimum node size used in the one-round monotone XGBoost calibrator.
    alpha : float, default=0.05
        Miscoverage level for confidence intervals.
    inference : {"wald", "bootstrap"}, default="wald"
        Type of confidence interval to report.
    n_bootstrap : int, default=999
        Number of bootstrap draws when inference="bootstrap".
    random_state : int, optional
        Seed for the bootstrap.
    max_depth : int, default=15
        Maximum depth for the calibration fit.

    Returns
    -------
    result : dict
        Dictionary containing the estimate, standard error, confidence
        interval, and the calibrated scores on both samples.

    Notes
    -----
    This implementation treats the initial score m(X) as fixed. If the score
    is itself estimated on the labeled data, cross-fitting should be used.
    """
    Y_labeled = _as_1d_array(Y_labeled, "Y_labeled")
    score_labeled, score_unlabeled = _resolve_scores(
        X_labeled=X_labeled,
        Y_labeled=Y_labeled,
        X_unlabeled=X_unlabeled,
        score_fn=score_fn,
        score_labeled=score_labeled,
        score_unlabeled=score_unlabeled,
    )

    n = len(Y_labeled)
    N = len(score_unlabeled)
    M = n + N
    rho_hat = n / M

    calibrator = isoreg_with_xgboost(
        x=score_labeled,
        y=Y_labeled,
        max_depth=max_depth,
        min_child_weight=min_obs,
    )

    calibrated_labeled = calibrator(score_labeled)
    calibrated_unlabeled = calibrator(score_unlabeled)

    psi_hat = (
        n * calibrated_labeled.mean() + N * calibrated_unlabeled.mean()
    ) / M

    # Estimated influence-function contributions for Wald inference.
    if_labeled = calibrated_labeled - psi_hat + (Y_labeled - calibrated_labeled) / rho_hat
    if_unlabeled = calibrated_unlabeled - psi_hat
    asymptotic_var = rho_hat * np.var(if_labeled, ddof=1) + (1 - rho_hat) * np.var(
        if_unlabeled, ddof=1
    )
    se_wald = np.sqrt(asymptotic_var / M)

    z = NormalDist().inv_cdf(1 - alpha / 2)
    ci_wald = (psi_hat - z * se_wald, psi_hat + z * se_wald)

    result = {
        "estimate": psi_hat,
        "standard_error": se_wald,
        "confidence_interval": ci_wald,
        "inference": "wald",
        "alpha": alpha,
        "rho_hat": rho_hat,
        "calibrator": calibrator,
        "score_labeled": score_labeled,
        "score_unlabeled": score_unlabeled,
        "calibrated_labeled": calibrated_labeled,
        "calibrated_unlabeled": calibrated_unlabeled,
    }

    if inference == "wald":
        return result

    if inference != "bootstrap":
        raise ValueError("inference must be either 'wald' or 'bootstrap'.")

    rng = np.random.default_rng(random_state)
    bootstrap_estimates = np.empty(n_bootstrap)

    for b in range(n_bootstrap):
        idx_l = rng.integers(0, n, size=n)
        idx_u = rng.integers(0, N, size=N)

        calibrator_b = isoreg_with_xgboost(
            x=score_labeled[idx_l],
            y=Y_labeled[idx_l],
            max_depth=max_depth,
            min_child_weight=min_obs,
        )
        calibrated_labeled_b = calibrator_b(score_labeled[idx_l])
        calibrated_unlabeled_b = calibrator_b(score_unlabeled[idx_u])

        bootstrap_estimates[b] = (
            n * calibrated_labeled_b.mean() + N * calibrated_unlabeled_b.mean()
        ) / M

    ci_boot = tuple(
        np.quantile(bootstrap_estimates, [alpha / 2, 1 - alpha / 2])
    )
    result.update(
        {
            "standard_error": np.std(bootstrap_estimates, ddof=1),
            "confidence_interval": ci_boot,
            "inference": "bootstrap",
            "bootstrap_estimates": bootstrap_estimates,
        }
    )
    return result


# Example 1: supply a score function m(X)
# fit = calibrated_plugin_ppi(
#     Y_labeled=Y,
#     X_labeled=X,
#     X_unlabeled=X_tilde,
#     score_fn=m,
#     min_obs=20,
#     inference="wald",
# )

# Example 2: supply precomputed scores m(X) and m(Xtilde)
# fit = calibrated_plugin_ppi(
#     Y_labeled=Y,
#     score_labeled=mX,
#     score_unlabeled=mX_tilde,
#     min_obs=20,
#     inference="bootstrap",
#     n_bootstrap=999,
#     random_state=1,
# )
\end{codeblock}

\subsection{Linear calibration}
\begin{codeblock}
import numpy as np

from statistics import NormalDist


def _as_1d_array(x, name):
    x = np.asarray(x).reshape(-1)
    if x.size == 0:
        raise ValueError(f"{name} must be non-empty.")
    return x


def _resolve_scores(
    X_labeled=None,
    Y_labeled=None,
    X_unlabeled=None,
    score_fn=None,
    score_labeled=None,
    score_unlabeled=None,
):
    if score_fn is not None:
        if X_labeled is None or X_unlabeled is None:
            raise ValueError(
                "If score_fn is provided, then X_labeled and X_unlabeled "
                "must also be provided."
            )
        score_labeled = score_fn(X_labeled)
        score_unlabeled = score_fn(X_unlabeled)
    elif score_labeled is None or score_unlabeled is None:
        raise ValueError(
            "Provide either score_fn with X_labeled/X_unlabeled, or provide "
            "both score_labeled and score_unlabeled."
        )

    score_labeled = _as_1d_array(score_labeled, "score_labeled")
    score_unlabeled = _as_1d_array(score_unlabeled, "score_unlabeled")

    if Y_labeled is not None and len(score_labeled) != len(Y_labeled):
        raise ValueError("score_labeled and Y_labeled must have the same length.")

    return score_labeled, score_unlabeled


def linear_calibrated_plugin_ppi(
    Y_labeled,
    X_labeled=None,
    X_unlabeled=None,
    score_fn=None,
    score_labeled=None,
    score_unlabeled=None,
    alpha=0.05,
):
    """
    Linear-calibrated plug-in PPI estimator for psi = E[Y].

    The calibration model is
        E[Y | m(X)] approx beta_0 + beta_1 m(X),
    estimated by least squares on the labeled sample. To avoid extrapolation,
    scores are clipped to the labeled score range before applying the
    calibrator to either sample.
    """
    Y_labeled = _as_1d_array(Y_labeled, "Y_labeled")
    score_labeled, score_unlabeled = _resolve_scores(
        X_labeled=X_labeled,
        Y_labeled=Y_labeled,
        X_unlabeled=X_unlabeled,
        score_fn=score_fn,
        score_labeled=score_labeled,
        score_unlabeled=score_unlabeled,
    )

    n = len(Y_labeled)
    N = len(score_unlabeled)
    M = n + N
    rho_hat = n / M

    design = np.column_stack([np.ones(n), score_labeled])
    beta_hat, _, _, _ = np.linalg.lstsq(design, Y_labeled, rcond=None)
    score_range = (score_labeled.min(), score_labeled.max())

    def calibrator(score):
        score = np.asarray(score).reshape(-1)
        score = np.clip(score, score_range[0], score_range[1])
        return beta_hat[0] + beta_hat[1] * score

    calibrated_labeled = calibrator(score_labeled)
    calibrated_unlabeled = calibrator(score_unlabeled)

    psi_hat = (
        n * calibrated_labeled.mean() + N * calibrated_unlabeled.mean()
    ) / M

    if_labeled = calibrated_labeled - psi_hat + (Y_labeled - calibrated_labeled) / rho_hat
    if_unlabeled = calibrated_unlabeled - psi_hat
    asymptotic_var = rho_hat * np.var(if_labeled, ddof=1) + (1 - rho_hat) * np.var(
        if_unlabeled, ddof=1
    )
    standard_error = np.sqrt(asymptotic_var / M)

    z = NormalDist().inv_cdf(1 - alpha / 2)
    confidence_interval = (
        psi_hat - z * standard_error,
        psi_hat + z * standard_error,
    )

    return {
        "estimate": psi_hat,
        "standard_error": standard_error,
        "confidence_interval": confidence_interval,
        "alpha": alpha,
        "rho_hat": rho_hat,
        "beta_hat": beta_hat,
        "score_range": score_range,
        "score_labeled": score_labeled,
        "score_unlabeled": score_unlabeled,
        "calibrated_labeled": calibrated_labeled,
        "calibrated_unlabeled": calibrated_unlabeled,
    }


# Example: linear calibration with a score function m(X)
# fit = linear_calibrated_plugin_ppi(
#     Y_labeled=Y,
#     X_labeled=X,
#     X_unlabeled=X_tilde,
#     score_fn=m,
# )
\end{codeblock}


\end{document}
